\subsection{Principal coverings}

DEFINITION 2. --- Let B be a topological space and let G be a group. Endow G with the discrete topology. A principal fiber bundle with base B and group G is called a principal covering of B with group G.

This terminology is legitimate because such a B-space is a covering of B.

A principal G-bundle over a nonempty topological space has a degree, equal to Card G.

PROPOSITION 4. — Let B be a topological space, let (E, p) be a B-space, and let G be a discrete topological group acting on the right on E. The following assertions are equivalent:

\begin{enumerate}
\def\labelenumi{(\roman{enumi})}
\item
  The B-space E is a principal covering with group G;
\item
  For every \(g \in G\) and every \(x \in E\) , we have \(p(x \cdot g) = p(x)\) , the map p induces a homeomorphism from E/G onto B, and the map \(\theta: (x, g) \mapsto (x, x \cdot g)\) is a homeomorphism from E × G onto \(E \times_{B} E\);
\item
  The group G acts continuously and freely on E, the map p is étale, and its fibers are the orbits of G.
\end{enumerate}

(i)⇒(ii): This follows from I, p. 91 and from proposition 3 of I, p. 94. (ii)⇒(iii): Under hypothesis (ii), the action law of G is continuous, and the map p is surjective and open (TG, III, p. 10, lemma 2). Let \(x \in E\) ; the map \(\theta\) induces a bijection from \(\{x\} \times G\) onto \(\{x\} \times p^{-1}(p(x))\) , so the group G acts freely and its orbits are the fibers of p. If e denotes the identity element of G, the diagonal of \(E \times_{B} E\) is the image of \(E \times \{e\}\) by the homeomorphism \(\theta\). Since the group G is discrete, the set \(\{e\}\) is open in G, hence the diagonal \(\Delta_{E}\) is open in \(E \times_{B} E\). By Proposition 7 of I, p. 31, the map p is étale.

(iii)⇒(i): Since every fiber of p is an orbit of the action of G on E, the map p is surjective. As p is étale, for every point b of B, there exists an open neighborhood U of b and a continuous section s of p over U. The set \(s(\mathrm{U})\) is open in E and s induces a homeomorphism from U onto \(s(\mathrm{U})\) (I, p. 30, cor. 3). For every \(g \in G\), the set \(s(\mathrm{U}) \cdot g\) is open in E and the union of these sets over all \(g \in G\) is equal to \(p^{-1}(\mathrm{U})\). If g and \(g'\) are two distinct elements of G, the sets \(s(\mathrm{U}) \cdot g\) and \(s(\mathrm{U}) \cdot g'\) are disjoint, since G acts freely on E. The map \(f: U \times G \to p^{-1}(\mathrm{U})\) defined by \(f(u, g) = s(u) \cdot g\) is continuous for \((u, g) \in U \times G\) and is therefore a homeomorphism from \(U \times G\) onto \(p^{-1}(\mathrm{U})\), compatible with the right action of G. By definition, E is therefore a principal covering of B with group G.

Examples. --- 1) Let E be a topological space, G a discrete topological group acting continuously and freely on E on the right. For the canonical map \(p \colon \mathrm{E} \to \mathrm{E} / \mathrm{G}\) to make E a principal covering of E/G, it is necessary and sufficient that it be étale (condition (iii) of Proposition 4). For example, suppose there exists a topological space X and an étale map \(q \colon \mathrm{E} \to \mathrm{X}\) compatible with the action of G on E; then E is a principal covering of E/G. Indeed, let us denote by \(q' \colon \mathrm{E} / \mathrm{G} \to \mathrm{X}\) the map induced by \(q\); we have \(q = q' \circ p\), and therefore \(p\) is étale by Proposition 6, d) of I, p. 29.

\begin{enumerate}
\def\labelenumi{\arabic{enumi})}
\setcounter{enumi}{1}
\tightlist
\item
  Let \(q \colon \mathrm{E} \to \mathrm{X}\) be a separated étale morphism. The group \(\operatorname{Aut}_{\mathrm{X}}(\mathrm{E})\) , endowed with the discrete topology, acts continuously on the left on E. If the space E is connected, this action is free (I, p. 34, corollary 2). Let G denote the group \(\operatorname{Aut}_{\mathrm{X}}(\mathrm{E})^{\circ}\) and endow E with the right action opposite to that of \(\operatorname{Aut}_{\mathrm{X}}(\mathrm{E})\) . By the preceding example, the space E, endowed with the canonical map \(p \colon \mathrm{E} \to \mathrm{E} / \mathrm{G}\) , is a principal covering with group G.
\end{enumerate}

\subsection{Proper and free actions of discrete groups}

THEOREM 1. --- Let G be a discrete group acting continuously on the right on a topological space E. Suppose that every point of E has a neighborhood U such that U ∩ (U · g) = ∅ for every element g of G other than the identity element. Then the canonical map p: E → E/G makes E a principal covering of E/G with group G.

The action of G on E is free by hypothesis, so it suffices to prove that the map p is étale (I, p. 98, example 1). Let x be a point of E and let U be an open neighborhood of x such that \(U \cap U \cdot g = \varnothing\) for every element g of G distinct from the identity element. The map p is open (TG, III, p. 10, lemma 2) and induces an injective, continuous, and open map from U onto \(p(\mathrm{U})\) , hence a homeomorphism, which proves that the map p is étale.

Examples. --- 1) Let n be an integer \(\geqslant 0\) , let \(\mathbf{P}_{n}(\mathbf{R})\) be the n-dimensional projective space (TG, VI, p. 13) and let \(S_{n}\) be the unit sphere in \(R^{n+1}\) (TG, VI, p. 9). The canonical map from \(S_{n}\) onto \(\mathbf{P}_{n}(\mathbf{R})\) makes \(S_{n}\) a principal covering with group \(\{1,-1\}\) , this group acting by homotheties; the orbits are the pairs of diametrically opposite points.

\begin{enumerate}
\def\labelenumi{\arabic{enumi})}
\setcounter{enumi}{1}
\item
  Every covering of degree 2 has a unique structure of a principal covering with group Z/2Z.
\item
  Let X be a finite-dimensional locally separated differentiable manifold over R, and \(\widetilde{X}\) be the manifold of orientations of the tangent bundle of X (VAR, R, 10.2.4). The space \(\widetilde{X}\) , equipped with its canonical projection onto X, is a principal covering of X with group \(\{1,-1\}\) .
\end{enumerate}

COROLLARY 1. --- Let G be a discrete topological group acting continuously on the right on a topological space E. The following conditions are equivalent:

\begin{enumerate}
\def\labelenumi{(\roman{enumi})}
\item
  The group G acts properly and freely on E;
\item
  The space E/G is Hausdorff, and E is a principal covering with base E/G and group G.
\end{enumerate}

Moreover, under these conditions, the space E is Hausdorff.

Suppose condition (i) is satisfied. Then the spaces E and E/G are separated (TG, III, p. 29, prop. 3), and for every \(x \in \mathrm{E}\) there exists an open neighborhood U of \(x\) in E such that \(\mathrm{U} \cap (\mathrm{U} \cdot g) = \varnothing\) for every element \(g\) of G other than the identity element (TG, III, p. 32, prop. 8). By th. 1, condition (ii) is then satisfied.

The implication (ii)⇒(i) follows from Corollary 1 of I, p. 96.

COROLLARY 2. --- Let G be a topological group and let H be a discrete subgroup of G. Let H act on G by right translations. Then the canonical map from G into the space G/H of left cosets modulo H endows G with a structure of a principal covering of G/H with group H.

Let V be a neighborhood of the identity element \(e\) of G such that \(\mathrm{H} \cap \mathrm{V} = \{e\}\) . There exists an open neighborhood U of \(e\) in G such that \(\mathrm{U}^{-1} \cdot \mathrm{U} \subset \mathrm{V}\) (TG, III, p. 3) and, consequently, \(\mathrm{U} \cap \mathrm{U} \cdot h = \varnothing\) for all \(h \in \mathrm{H}\) , \(h \neq e\) . Corollary 2 of I, p. 100 therefore follows from Theorem 1.

Example 4. --- Equipped with the canonical map from R to T = R/Z (TG, V, p. 2), R is a principal covering of R/Z with group Z.

Remark 1. --- Let G be a separated topological group, K a compact subgroup of G, and H a discrete subgroup of G. The group G acts continuously and properly on the right on the space K\G of right cosets modulo K (TG, III, p. 30, corollary). Since the group H is closed in G, it also acts properly on K\G (TG, III, p. 27, example 1). If, moreover, one has \(H \cap gKg^{-1} = \{e\}\) for all \(g \in G\), the group H acts freely on K\G. The space K\G is then a principal covering with group H of K\G/H by Corollary 1.

COROLLARY 3. --- Let G and G' be topological groups, and let \(\varphi\colon G \to G'\) be a continuous, open, and surjective homomorphism. Suppose that the kernel H of \(\varphi\) is discrete. Then, for the action of H on G by right translations, \(\varphi\) makes G a principal covering of G' with group H.

If moreover the group G is connected, then H is contained in the center of G.

The homomorphism \(\varphi\) induces an isomorphism of topological groups from G/H onto G' (TG, III, p. 16, prop. 24), whence the first assertion by Corollary 2. Suppose the group G is connected. For every \(h \in \mathrm{H}\) , the continuous map from G into H defined by \(g \mapsto ghg^{-1}\) is constant, with value \(h\) . The group H is therefore contained in the center of G.

Remarks. --- 2) Let G and G' be topological groups, and let \(\varphi\colon\mathrm{G}\to\mathrm{G}'\) be a continuous and open homomorphism. If the group G' is connected, the homomorphism \(\varphi\) is surjective. Indeed, \(\varphi(\mathrm{G})\) is an open, hence closed, subgroup of \(\mathrm{G}'\), and consequently is equal to \(\mathrm{G}'\).

\begin{enumerate}
\def\labelenumi{\arabic{enumi})}
\setcounter{enumi}{2}
\tightlist
\item
  Let G be a locally compact group which is countable at infinity, and let G' be a separated topological group whose underlying space is a Baire space. Every continuous surjective homomorphism from G onto G' is open (TG, IX, p. 56, corollary and TG, III, p. 16, prop. 24).
\end{enumerate}

Examples. --- 5) For every integer \(n > 0\), let us denote by \(\mu_{n}\) the group of n-th roots of unity in C (A, V, p. 75). The map \(z \mapsto z^{n}\) makes \(C^{*}\) a principal covering of \(C^{*}\) with group \(\mu_{n}\), this group acting in \(C^{*}\) by multiplication.

\begin{enumerate}
\def\labelenumi{\arabic{enumi})}
\setcounter{enumi}{5}
\tightlist
\item
  Equipped with the map \(z \mapsto e^{2\pi iz}\) (TG, VIII, p. 8, remark), the space C is a principal covering of \(C^{*}\) with group Z. The map \(x \mapsto e^{2\pi ix} = \mathbf{e}(x)\) from R to \(S_{1}\) makes R a principal covering of \(S_{1}\) with group Z.
\end{enumerate}

\subsection{Galois coverings}

DEFINITION 3. --- Let B be a nonempty topological space. A covering E of B is said to be Galois if it is connected and if, for every point b of B, the action of the group \(\mathrm{Aut}_{B}(E)\) of B-automorphisms of E on the fiber \(E_{b}\) is transitive.

Let E be a Galois covering of a topological space B and let p be its projection. The map p is surjective (A, I, p. 56, def. 6), so the space B is connected. Consequently, the covering (E, p) has a degree, which is nonzero.

PROPOSITION 5. --- Let B be a nonempty topological space and let E be a Galois covering of B. The map \((h,x)\mapsto h(x)\) from \(\mathrm{Aut}_{\mathrm{B}}(\mathrm{E})\times\mathrm{E}\) into E is a right action of the group \(\mathrm{Aut}_{\mathrm{B}}(\mathrm{E})^{\circ}\) on E which makes E a principal covering with group \(\mathrm{Aut}_{\mathrm{B}}(\mathrm{E})^{\circ}\) .

Endow the group \(\mathrm{Aut}_{\mathrm{B}}(\mathrm{E})^{\circ}\) with the discrete topology; the operation law \((h,x)\mapsto h(x)\) is then continuous. This action is free (I, p. 34, corollary 2 of proposition 11). Since E is a Galois covering of B, its fibers are the orbits of this action. By proposition 4 of I, p. 97, E is a principal covering with group \(\mathrm{Aut}_{\mathrm{B}}(\mathrm{E})^{\circ}\) .

THEOREM 2. --- Let B be a connected topological space, and let E be a connected, nonempty covering of B. The following properties are equivalent:

\begin{enumerate}
\def\labelenumi{(\roman{enumi})}
\item
  The covering E is Galois.
\item
  There exists a discrete topological group G and a continuous right action of G on E that makes E a principal covering with group G.
\item
  The covering \(E \times_{B} E\) of E defined by the projection \(\mathrm{pr}_{1}\) is trivializable.
\end{enumerate}

When these conditions are satisfied, the canonical map \(G \to \mathrm{Aut}_{B}(E)^{\circ}\) defined by the action of G is an isomorphism of groups.

If, moreover, the space B is locally connected, the preceding properties are equivalent to the following property:

(i') There exists a point b of B such that the action of the group \(\mathrm{Aut}_{\mathrm{B}}(\mathrm{E})\) on the fiber \(E_{b}\) is transitive.

The implication (i)⇒(ii) follows from Proposition 5, and the implication (ii)⇒(iii) from Remark 1 of I, p. 95. Let us prove (iii)⇒(i). Denote by \(p\colon E \to B\) the projection of the covering E. Since B is connected, this covering has a degree, and this degree is nonzero because E is not empty. Let \(b\) be a point of B. Let us prove that the action of \(\mathrm{Aut}_{B}(E)\) on the fiber \(E_{b}\) is transitive. By the preceding, this fiber is nonempty. Let \(x\) and \(x'\) be points of \(E_{b}\). The B-morphisms \(h\colon E \to E\) such that \(h(x) = x'\) correspond bijectively to continuous sections \(s\) of the map \(\mathrm{pr}_{1}\colon E \times_{B} E \to E\) such that \(s(x) = (x, x')\) (I, p. 9, prop. 3). Under hypothesis (iii), such a section exists (I, p. 70, cor. 2) and is unique (I, p. 34, cor. 1 of prop. 11) because the space E is connected. Thus, for every pair \((x, x') \in E \times_{B} E\) there exists a unique B-morphism \(h\colon E \to E\) such that \(h(x) = x'\) . If \(h'\colon E \to E\) is the unique B-morphism such that \(h'(x') = x\) , we have \(h'(h(x)) = x\) and \(h(h'(x')) = x'\) , whence \(h' \circ h = \text{Id}_E\) and \(h \circ h' = \text{Id}_E\) . This proves that \(h\) is a B-automorphism of E. The covering E is therefore Galois.

We have shown that conditions (i), (ii), (iii) are equivalent. Suppose they are satisfied. Let \(\delta\colon\mathbf{G}\to\mathrm{Aut}_{\mathrm{B}}(\mathrm{E})\) be the map which associates with \(g\in\mathrm{G}\) the B-automorphism \(x\mapsto x\cdot g\) of E. The map \(\delta\) is a group homomorphism from G into \(\mathrm{Aut}_{\mathrm{B}}(\mathrm{E})^{\circ}\). Since G acts freely on E, the map \(\delta\) is injective. Let \(h\in\mathrm{Aut}_{\mathrm{B}}(\mathrm{E})\), let \(x\in\mathrm{E}\), and let \(g\) be the unique element of G such that \(x\cdot g=h(x)\).

B-morphisms \(h\) and \(\delta(g)\) coincide at \(x\) and therefore everywhere, since the space E is connected (I, p. 34, cor. 1 of prop. 11), which proves that \(\delta\) is an isomorphism.

Suppose now that the space B is locally connected, and let us prove, under hypothesis (i'), that the covering E is Galois. Let \(E'\) be the quotient space of E by the right action of \(\mathrm{Aut}_{B}(E)^{\circ}\) and denote by \(q\colon E \to E'\) the canonical map. It makes E a surjective covering of E' (I, p. 98, example 2); the map \(p\colon E \to B\) passes to the quotient to define a continuous map \(p'\colon E' \to B\) such that \(p'\circ q = p\) . Since the space B is locally connected, the B-space \((E', p')\) is a covering (I, p. 81, prop. 7). Under hypothesis (i'), there exists a point \(b\) of B such that the fiber \(E'_{b}\) has exactly one element; since the space B is connected, the same holds for every point \(b\) of B (I, p. 74, prop. 4), which proves that E is a Galois covering of B.

PROPOSITION 6. --- Let B be a topological space and G a group. Let (E, p) be a principal covering of B with group G. Suppose the space B is connected and locally connected. Let E₀ be a connected component of E and let G₀ be the subgroup of G stabilizing E₀ (A, I, p. 51). The B-space (E₀, p \textbar{} E₀) is a principal covering for the right action of G₀ on E₀. This covering is Galois.

Since the space B is locally connected, so is E, so that \(E_{0}\) is an open subspace of E (TG, I, p. 85, prop. 11). Since it is also closed (TG, I, p. 83), the space \(E_{0}\) is therefore a covering of B (I, p. 80, cor. 1). Since B is connected, this covering has a degree; since \(E_{0}\) is nonempty, all fibers of \(p|E_{0}\) are therefore nonempty. Let x be a point of \(E_{0}\), and let \(x' \in E_{0}\) such that \(p(x') = p(x)\); there exists an element \(g \in G\) such that \(x \cdot g = x'\). Since the connected components \(E_{0}\) and \(E_{0} \cdot g\) have a common point, they are equal, which implies \(g \in G_{0}\). Thus, the group \(G_{0}\) acts transitively on the fiber of the B-space \(E_{0}\) at \(p(x)\). Proposition 6 follows.

Note. --- If, in Proposition 6, one does not assume the space B to be locally connected, it may happen that the space \(E_{0}\) is not a covering of B (I, p. 147, exerc. 1).

\subsection{Associated fiber bundles}

Let E and F be sets, and let G be a group acting on the right on E and on the left on F. The group G acts on the right on the product \(E \times F\) by the law \((x, y) \cdot g = (x \cdot g, g^{-1} \cdot y)\) , for \(g \in G\) and \((x, y) \in E \times F\) . The quotient set of \(E \times F\) by this operation is denoted \(E \times^{G} F\) .

When E and F are topological spaces, one equips the set \(E \times^{G} F\) with the quotient topology induced by that of \(E \times F\) . The canonical map of \(E \times F\) onto \(E \times^{G} F\) is continuous. If the group G acts continuously on E and F, it is open.

Additionally, let \(F'\) be a set on which G acts on the left, and let \(h: F \to F'\) be a map compatible with the operations of G on F and \(F'\) (A, I, p. 50). The map \(Id_{E} \times h: E \times F \to E \times F'\) is compatible with the operations of G and, by passing to quotients, defines a map denoted \(Id_{E} \times^{G} h\) from \(E \times^{G} F\) into \(E \times^{G} F'\) .

If \(h \colon F \to F'\) is a continuous map (compatible with the operations of G on F and \(F'\) ), the map \(\mathrm{Id}_{\mathrm{E}} \times^{\mathrm{G}} h\) is continuous (TG, I, p. 21, prop. 6).

Examples. --- 1) Let F be a topological space and let G be a topological group acting continuously on the left on F. If one endows the topological space G with the action of G by right translations, the space G × \(^{G}\) F is canonically identified with F as follows. The continuous maps \(\varphi\) : F → G × F and \(\psi\) : G × F → F defined by \(\varphi(f) = (e, f)\) (where e denotes the identity element of G) and \(\psi(g, f) = g \cdot f\) induce continuous maps \(\overline{\varphi}\) : F → G × \(^{G}\) F and \(\overline{\psi}\) : G × \(^{G}\) F → F which are inverses of each other.

\begin{enumerate}
\def\labelenumi{\arabic{enumi})}
\setcounter{enumi}{1}
\item
  Example 1 generalizes as follows. Let B and F be topological spaces and let G be a topological group acting continuously on the left on F. By passing to quotients, the map from B × F to (B × G) × F given by (b, f) ↦ ((b, e), f) and the map from (B × G) × F to B × F given by ((b, g), f) ↦ (b, gf) define B-isomorphisms B × F → (B × G) × \(^{G}\) F and (B × G) × \(^{G}\) F → B × F which are inverses of each other.
\item
  Similarly, let E be a topological space and let G be a topological group acting continuously on the right on E. If the topological space G is endowed with the action of G by left translations, the space \(E \times^{G} G\) is identified with E.
\end{enumerate}

Let E and F be topological spaces, and let G be a topological group acting continuously on the right on E and on the left on F. Let B be a topological space and \(p\colon E \to B\) a continuous map such that \(p(x \cdot g) = p(x)\) for \(x \in E\) and \(g \in G\) . The map \(p \circ pr_{1}\colon E \times F \to B\) defines, by passing to the quotient, a continuous map \(p^{F}\colon E \times^{G} F \to B\) and the canonical map \(\pi\colon E \times F \to E \times^{G} F\) is a B-morphism.

Let B' be a topological space and let \(h\colon B'\to B\) be a continuous map. The group G acts continuously on the right on \(B'\times_{B}E\) by the operation law \(((b',x),g)\mapsto(b',x\cdot g)\). By composing the canonical isomorphism \(((b',x),y)\mapsto(b',(x,y))\) of \((\mathrm{B}'\times_{\mathrm{B}}\mathrm{E})\times\mathrm{F}\) onto \(\mathrm{B}'\times_{\mathrm{B}}(\mathrm{E}\times\mathrm{F})\) and the map \(\operatorname{Id}_{\mathrm{B}'}\times\pi\) from \(\mathrm{B}'\times_{\mathrm{B}}(\mathrm{E}\times\mathrm{F})\) into \(\mathrm{B}'\times_{\mathrm{B}}(\mathrm{E}\times^{\mathrm{G}}\mathrm{F})\), we obtain a continuous map \(\lambda_{0}\colon(\mathrm{B}'\times_{\mathrm{B}}\mathrm{E})\times\mathrm{F}\to\mathrm{B}'\times_{\mathrm{B}}(\mathrm{E}\times^{\mathrm{G}}\mathrm{F})\). By passing to the quotient, it defines a continuous map

\[
\lambda \colon (\mathrm{B} ^ {\prime} \times_ {\mathrm{B}} \mathrm{E}) \times^ {\mathrm{G}} \mathrm{F} \rightarrow \mathrm{B} ^ {\prime} \times_ {\mathrm{B}} (\mathrm{E} \times^ {\mathrm{G}} \mathrm{F}).
\]

Lemma. --- The map λ is a homeomorphism.

The map \(\lambda_0\) is surjective, and two elements of \((\mathrm{B}' \times_{\mathrm{B}} \mathrm{E}) \times \mathrm{F}\) have the same image under \(\lambda_0\) if and only if they belong to the same orbit for the action of G on \((\mathrm{B}' \times_{\mathrm{B}} \mathrm{E}) \times \mathrm{F}\). It follows that the map \(\lambda\) is bijective. Since the map \(\pi\) is open, so is the map \(\mathrm{Id}_{\mathrm{B}'} \times_{\mathrm{B}} \pi\) (I, p. 17, prop. 8), hence so are \(\lambda_0\) and \(\lambda\) (TG, I, p. 32, prop. 3), which proves that \(\lambda\) is a homeomorphism.

PROPOSITION 7. --- Let B be a topological space, let G be a topological group, let (E,p) be a principal fiber bundle over B with group G, and let F be a topological space on which G acts continuously on the left.

\begin{enumerate}
\def\labelenumi{\alph{enumi})}
\item
  The topological space \(E \times^{G} F\) together with the continuous map \(p^{F}: E \times^{G} F \to B\) deduced from the map \(p \circ pr_{1}: E \times F \to B\) by passing to the quotient is a locally trivial fiber bundle of fiber type F; it is trivializable if the B-bundle E is trivializable.
\item
  Let \(\pi\colon\mathrm{E}\times\mathrm{F}\to\mathrm{E}\times^{\mathrm{G}}\mathrm{F}\) be the canonical surjection. The map \(\mu\colon\mathrm{E}\times\mathrm{F}\to\mathrm{E}\times_{\mathrm{B}}(\mathrm{E}\times^{\mathrm{G}}\mathrm{F})\) which associates \((x,\pi(x,f))\) with \((x,f)\) is a homeomorphism whose inverse map is a trivialization of the locally trivial fiber bundle over E \((\mathrm{E}\times_{\mathrm{B}}(\mathrm{E}\times^{\mathrm{G}}\mathrm{F}),\mathrm{pr}_{1})\).
\end{enumerate}

Suppose first that E is the trivial principal fiber bundle \(B \times G\) . By Example 2, the space \(E \times^{G} F\) is then identified with \(B \times F\) and the map \(p^{F}\) with the first projection \(pr_{1}: B \times F \to B\) , which proves that \((\mathrm{E} \times^{\mathrm{G}}\mathrm{F}, p^{\mathrm{F}})\) is a trivializable fiber bundle over B in this case. In the general case, every point of B has a neighborhood U such that the map \(p_{\mathrm{U}}: p^{-1}(\mathrm{U}) \to \mathrm{U}\) makes \(p^{-1}(\mathrm{U})\) a trivializable principal fiber bundle over U. By the preceding, \((p^{-1}(\mathrm{U}) \times^{\mathrm{G}}\mathrm{F} \to \mathrm{U}, (p_{\mathrm{U}})^{\mathrm{F}})\) is a trivializable fiber bundle over U. By the lemma above, applied to the case where \(B' = U\) and \(h: U \to B\) is the canonical injection, the same holds for the U-space \(((p^{\mathrm{F}})^{-1}(\mathrm{U}), (p^{\mathrm{F}})_{\mathrm{U}})\) deduced from \((\mathrm{E} \times^{\mathrm{G}}\mathrm{F}, p^{\mathrm{F}})\) by restriction over U. This proves that \((\mathrm{E} \times^{\mathrm{G}}\mathrm{F}, p^{\mathrm{F}})\) is a locally trivial fiber bundle over B and concludes the proof of assertion a).

The map \(\theta\colon\mathrm{E}\times\mathrm{G}\to\mathrm{E}\times_{\mathrm{B}}\mathrm{E}\) defined by \(\theta(x,g)=(x,x\cdot g)\) is a homeomorphism (I, p. 94, prop. 3) compatible with the operations of G on E \(\times\) G and on E \(\times_{\mathrm{B}}\) E given by \(((x,g),g')\mapsto(x,gg')\) and \(((x,y),g')\mapsto(x,yg')\) respectively. By passing to quotients, \(\theta\) therefore induces a homeomorphism \(\theta'\) from (E \(\times\) G) \(\times^{\mathrm{G}}\) F onto (E \(\times_{\mathrm{B}}\) E) \(\times^{\mathrm{G}}\) F. The map \(\mu\) is the composition of the homeomorphism E \(\times\) F \(\rightarrow\) (E \(\times\) G) \(\times^{\mathrm{G}}\) F (example 2), of \(\theta'\) and of the canonical homeomorphism (E \(\times_{\mathrm{B}}\) E) \(\times^{\mathrm{G}}\) F \(\rightarrow\) E \(\times_{\mathrm{B}}\) (E \(\times^{\mathrm{G}}\) F) (I, p. 105, lemma), whence \(b\) ).

Under the hypotheses of Proposition 7, the locally trivial fiber bundle (E × \(^{G}\) F, p \(^{F}\) ) is called the locally trivial fiber bundle of fiber type F associated with the principal fiber bundle (E, p). All fibers of the B-space E × \(^{G}\) F are homeomorphic to the space F. In particular, if the space F is discrete, the map p \(^{F}\) is a covering.

Examples. --- 4) Let B be a topological space, G a topological group, and let (E, p) be a principal G-bundle over B. Let H be a subgroup of G.

Denote by \(\varphi\colon\mathrm{E}\to\mathrm{E}\times(\mathrm{G}/\mathrm{H})\) and \(\psi\colon\mathrm{E}\times(\mathrm{G}/\mathrm{H})\to\mathrm{E}/\mathrm{H}\) the maps defined by \(\varphi(x)=(x,\mathrm{H})\) and \(\psi(x,g\mathrm{H})=(x\cdot g)\mathrm{H}\) . They are compatible with the projections to B and with the operations of G, and they define, by passing to quotients, morphisms of B-spaces \(\overline{\varphi}\colon\mathrm{E}/\mathrm{H}\to\mathrm{E}\times^{\mathrm{G}}(\mathrm{G}/\mathrm{H})\) and \(\overline{\psi}\colon\mathrm{E}\times^{\mathrm{G}}(\mathrm{G}/\mathrm{H})\to\mathrm{E}/\mathrm{H}\) , inverses of each other. We say that \(\overline{\varphi}\) is the canonical homeomorphism from E/H onto \(\mathrm{E}\times^{\mathrm{G}}(\mathrm{G}/\mathrm{H})\) . In particular, the topological space E/H equipped with the continuous map \(p_{\mathrm{H}}\colon\mathrm{E}/\mathrm{H}\to\mathrm{B}\) is a locally trivial fiber bundle over B with fiber type G/H.

If moreover H is a normal subgroup of G, the action of G endows, by passing to quotients, the B-space E/H with the structure of a principal fiber bundle with group G/H.

In particular, if E is a principal covering of B with group G, then E/H is a covering of B; if H is normal in G, it is a principal covering with group G/H.

\begin{enumerate}
\def\labelenumi{\arabic{enumi})}
\setcounter{enumi}{4}
\item
  Let B be a topological space, let G be a topological group, and let E be a principal G-bundle over B. Let F be a topological homogeneous space with respect to G (TG, III, p. 12). Let y be a point of F, and let \(G_{y}\) be its stabilizer. The map \(\varphi_{y}\colon x\mapsto(x,y)\) from E to \(E\times F\) defines, by passing to the quotient, a homeomorphism \(\overline{\varphi}_{y}\) of \(E/G_{y}\) onto \(E\times^{G}F\) . When the group G is abelian, the subgroup \(G_{y}\) does not depend on the point y, but the homeomorphism \(\overline{\varphi}_{y}\) in general depends on it.
\item
  Let B be a topological space, let G be a topological group, let \((E,p)\) be a principal G-bundle over B. Let H be a topological group and let \(f:G\to H\) be a morphism of topological groups; endow H with the left action of G given by \(g\cdot h=f(g)h\) .
\end{enumerate}

Let \(q \colon \mathrm{E} \times \mathrm{H} \to \mathrm{E} \times^{\mathrm{G}} \mathrm{H}\) be the canonical map; it is open, hence universally strict (I, p. 20, corollary 11). The map \(m \colon (x, h, h') \mapsto q(x, hh')\) from \(\mathrm{E} \times \mathrm{H} \times \mathrm{H}\) into \(\mathrm{E} \times^{\mathrm{G}} \mathrm{H}\) is continuous. Let \((x, h) \in \mathrm{E} \times \mathrm{H}\) , \(h' \in \mathrm{H}\) and \(g \in \mathrm{G}\) ; we have \(m(xg, f(g)^{-1}h, h') = q(xg, f(g)^{-1}hh') = q(x, hh') = m(x, h, h')\) . Consequently, there exists a unique map

\[
m ^ {\prime} \colon (\mathrm{E} \times^ {\mathrm{G}} \mathrm{H}) \times \mathrm{H} \rightarrow \mathrm{E} \times^ {\mathrm{G}} \mathrm{H}
\]

such that \(m'(q(x,h),h') = q(x,hh')\) for all \(x \in E\) and all \(h, h' \in H\) . Since q is universally strict, the map \(m'\) is continuous. This is a right action of the topological group H on the B-space \(E \times^{G} H\) .

Let U be an open subset of B such that \(E_{U}\) is isomorphic to the principal fiber bundle over U \(U \times G\) . Equipped with the operation of H, the U-space \((\mathrm{E} \times^{\mathrm{G}}\mathrm{H})_{\mathrm{U}}\) is identified with the principal fiber bundle \(U \times H\) . This proves that \(E \times^{G}H\) is a principal H-bundle over B.

DEFINITION 4. --- Let B be a topological space and let G be a topological group. A locally trivial fiber bundle X over B is said to be associated with a principal fiber bundle E over B with group G if there exists a topological space F on which the group G acts continuously on the left and a B-isomorphism from E × \(^{G}\) F onto X.

Let E be a principal G-bundle over B, and let X be a locally trivial fiber bundle over B associated with E. If the principal bundle E is trivializable, then X is trivializable (prop. 7). If B' is a topological space and \(h \colon \mathrm{B}' \to \mathrm{B}\) is a continuous map, the locally trivial fiber bundle \(\mathrm{B}' \times_{\mathrm{B}} \mathrm{X}\) deduced from X by base change is associated to the B'-principal fiber bundle \(\mathrm{B}' \times_{\mathrm{B}} \mathrm{E}\) (I, p. 105, lemma).

\subsection{Associated coverings}

Let B be a topological space, let G be a discrete topological group, and let E be a principal covering of B with group G. Let F be a G-set; if F is endowed with the discrete topology, the group G acts continuously on F. The B-space \(E \times^{G} F\) is then a covering. It is called the covering of B with fiber type F associated with the principal covering E.

DEFINITION 5. --- Let B be a topological space, let G be a discrete topological group, and let E be a principal covering of B with group G. A covering X of B is said to be associable to the principal covering E if there exists a G-set F and a B-isomorphism from E × \(^{G}\) F onto X.

Let B be a topological space, let G be a discrete topological group, and let E be a principal covering of B with group G.

For every covering X of B, the group G acts on the left on \(\mathcal{C}_{\mathrm{B}}(\mathrm{E};\mathrm{X})\) by the law defined by \((g\cdot h)(x)=h(x\cdot g)\) for \(x\in E\) , \(g\in G\) , \(h\in\mathcal{C}_{\mathrm{B}}(\mathrm{E};\mathrm{X})\) .

For every G-set F equipped with the discrete topology, define a map \(\alpha_{\mathrm{F}}\colon \mathrm{F}\to \mathcal{C}_{\mathrm{B}}(\mathrm{E};\mathrm{E}\times^{\mathrm{G}}\mathrm{F})\) by:

\[
\alpha_ {\mathrm{F}} (y) (x) = \pi (x, y) \quad \text {for } y \in \mathrm{F} \text { and } x \in \mathrm{E},\tag{2}
\]

where \(\pi\colon\mathrm{E}\times\mathrm{F}\to\mathrm{E}\times^{\mathrm{G}}\mathrm{F}\) is the canonical surjection. The map \(\alpha_{\mathrm{F}}\) is compatible with the operations of G on F and on \(\mathcal{C}_{\mathrm{B}}(\mathrm{E};\mathrm{E}\times^{\mathrm{G}}\mathrm{F})\) .

For every covering X of B, there exists a unique map \(\beta_{X}\) from \(E \times^{G} \mathcal{C}_{B}(E; X)\) into X such that:

\[
\beta_ {\mathrm{X}} (\pi (x, h)) = h (x) \quad \text {for } h \in \mathcal {C} _ {\mathrm{B}} (\mathrm{E}; \mathrm{X}) \text { and } x \in \mathrm{E}.\tag{3}
\]

Indeed, we have \((g^{-1}\cdot h)(x\cdot g)=h(x)\) for \(x\in E\) , \(g\in G\) and \(h\in\mathcal{C}_{\mathrm{B}}(\mathrm{E};\mathrm{X})\) , by definition of the action of G on \(\mathcal{C}_{\mathrm{B}}(\mathrm{E};\mathrm{X})\). When one equips the set \(\mathcal{C}_{\mathrm{B}}(\mathrm{E};\mathrm{X})\) with the discrete topology, the map \(\beta_{X}\) is a B-morphism of coverings.

When \(\mathrm{X} = \mathrm{E}\times^{\mathrm{G}}\mathrm{F}\) , we have

\[
\beta_ {\mathrm{X}} \left(\pi \left(x, \alpha_ {\mathrm{F}} (y)\right)\right) = \pi (x, y) \quad \text {   for   } x \in \mathrm{X} \text {   and   } y \in \mathrm{F}.\tag{4}
\]

PROPOSITION 8. --- Let B be a connected and nonempty topological space. Let G be a discrete group and let (E, p) be a principal covering of B with group G. Suppose that E is connected. With the notations above, we have:

\begin{enumerate}
\def\labelenumi{\alph{enumi})}
\item
  For every G-set F equipped with the discrete topology, the map \(\alpha_{F}\) is an isomorphism of the G-set F onto the G-set \(\mathcal{C}_{\mathrm{B}}(\mathrm{E};\mathrm{E}\times^{\mathrm{G}}\mathrm{F})\) .
\item
  Let X be a covering of B. The B-morphism \(\beta_{X}\) is an isomorphism from \(E \times^{G} \mathcal{C}_{B}(E; X)\) onto X if and only if the covering \((E \times_{B} X, pr_{1})\) of E is trivializable.
\item
  Let \(y\) and \(y'\) be points of F such that \(\alpha_{\mathrm{F}}(y) = \alpha_{\mathrm{F}}(y')\). The space E is not empty; let us choose a point \(x\) in it. We have \(\pi(x, y) = \pi(x, y')\) in \(\mathrm{E} \times^{\mathrm{G}} \mathrm{F}\). Consequently, there exists \(g \in \mathrm{G}\) such that \(x \cdot g = x\) and \(g^{-1} \cdot y = y'\). The first equality implies that \(g\) is the identity element \(e\) of G, hence \(y = y'\). Thus, the map \(\alpha_{\mathrm{F}}\) is injective. On the other hand, let \(h \in \mathcal{C}_{\mathrm{B}}(\mathrm{E}; \mathrm{E} \times^{\mathrm{G}} \mathrm{F})\) and let \(x\) be a point of E. Let \(x' \in \mathrm{E}\) and \(y' \in \mathrm{F}\) be such that \(h(x) = \pi(x', y')\). In particular, \(p(x) = p(x')\); there then exists an element \(g\) of G such that \(x' = x \cdot g\), and one also has \(h(x) = \pi(x, y)\), where we have set \(y = g \cdot y'\). The B-morphisms \(h\) and \(\alpha_{\mathrm{F}}(y)\) coincide at the point \(x\) of E; they are therefore equal since the space E is connected (I, p. 34, cor. 1 of prop. 11), and this proves that the map \(\alpha_{\mathrm{F}}\) is surjective.
\item
  By proposition 7, b) of I, p. 105, applied to F = \(\mathcal{C}_{\mathrm{B}}(\mathrm{E};\mathrm{X})\), the covering \(p^{*}(\mathrm{E}\times^{\mathrm{G}}\mathcal{C}_{\mathrm{B}}(\mathrm{E};\mathrm{X}))\) of E is isomorphic to the trivial covering \(\mathrm{E}\times \mathcal{C}_{\mathrm{B}}(\mathrm{E};\mathrm{X})\). If \(\beta_{\mathrm{X}}\) is an isomorphism, the covering \(p^{*}(\mathrm{X})\) of E is therefore trivializable. Conversely, suppose the covering \(p^{*}(\mathrm{X})\) is trivializable and let us show that the B-morphism \(\beta_{\mathrm{X}}\) is bijective; it will follow that \(\beta_{\mathrm{X}}\) is a B-isomorphism (I, p. 30, cor. 2 of prop. 6).
\end{enumerate}

The map \(\beta_{\mathrm{X}}\) is obtained by passing to the quotient from the map \(\gamma\colon\mathrm{E}\times\mathcal{C}_{\mathrm{B}}(\mathrm{E};\mathrm{X})\to\mathrm{X}\) defined by \(\gamma(x,h)=h(x)\) . Let us show that the map \(\gamma\) is surjective. Let \(x\) be a point of X. The projection of the B-space E is surjective, since it is a principal covering; there therefore exists a point \(y\) of E such that \((y,x)\in\mathrm{E}\times_{\mathrm{B}}\mathrm{X}\) . There then exists a continuous section \(s\) of the trivializable covering \(\mathrm{pr}_{1}:\mathrm{E}\times_{\mathrm{B}}\mathrm{X}\to\mathrm{E}\) such that \(s(y)=(y,x)\) . The map \(h=\mathrm{pr}_{2}\circ s:\mathrm{E}\to\mathrm{X}\) is a B-morphism and one has \(h(y)=x\) , whence the surjectivity of the map \(\gamma\) and, consequently, that of the map \(\beta_{\mathrm{X}}\) .

Let us finally show that \(\beta_{\mathrm{X}}\) is injective. Let \((x,h)\) and \((x^{\prime},h^{\prime})\) be elements of \(\mathrm{E}\times \mathcal{C}_{\mathrm{B}}(\mathrm{E};\mathrm{X})\) such that \(h(x) = h'(x')\) . Note that \(x\) and \(x'\) have the same projection in B; there therefore exists an element \(g\) of G such that \(x' = x\cdot g\) . One then has \(h(x) = h'(x\cdot g) = (g\cdot h')(x)\) . Since the space E is connected, one has \(h = g\cdot h'\) (I, p. 34, cor. 1 of prop. 11). Thus, \((x,h)\) and \((x',h')\) have the same class in \(\mathrm{E}\times^{\mathrm{G}}\mathcal{C}_{\mathrm{B}}(\mathrm{E};\mathrm{X})\) , which shows that the map \(\beta_{\mathrm{X}}\) is injective, which completes the proof.

COROLLARY 1. --- Let (E,p) be a principal covering of B with group G; suppose that E is connected and nonempty. A covering X of B is associable to E if and only if the covering \(p^{*}(X)\) is trivializable.

If the covering \(p^{*}(\mathrm{X})\) is trivializable, it follows from Proposition 8, b), that the covering X is associable to E. In this case, the map \(\beta_{X}\) identifies the covering X with the fiber-type covering \(\mathcal{C}_{\mathrm{B}}(\mathrm{E};\mathrm{X})\) associated with E. Conversely, let F be a G-set equipped with the discrete topology, and suppose that we have \(X = E \times^{G} F\) . Then, \(\alpha_{F}\) is an isomorphism (loc. cit., a)) and formula (4) implies that \(\beta_{X}\) is an isomorphism. Consequently, the covering \(p^{*}(\mathrm{X})\) is trivializable (prop. 8, b)), whence the corollary.

COROLLARY 2. --- Let B be a connected and locally connected topological space, let (E,p) be a principal covering of B with group G. Let E₀ be a connected component of E and let G₀ be the subgroup of G stabilizing E₀. The B-space (E₀, p \textbar{} E₀) is a principal covering with group G₀ (I, p. 103, prop. 6).

Every covering X of B that is associable to the principal covering E is associable to the principal covering \(E_{0}\). In particular, the covering E is associated with the principal covering \(E_{0}\).

Indeed, if the covering X is associable to E, the covering \(p^{*}(\mathrm{X})\) is trivializable and the covering \(p_{0}^{*}(\mathrm{X})\) induced by \(p^{*}(\mathrm{X})\) over \(E_{0}\) is therefore trivializable.

More precisely, note that the map \((x,g)\mapsto x\cdot g\) from \(E_{0}\times G\) into E induces, by passing to the quotient, a B-isomorphism of principal coverings from \(E_{0}\times^{G_{0}}G\) onto E.

PROPOSITION 9. --- Let B be a topological space, let E be a principal covering of B with group G, connected and nonempty. Let F be a nonempty G-set equipped with the discrete topology. For the space E × \(^{G}\) F to be connected, it is necessary and sufficient that G acts transitively on F.

Let U be an open and closed subset of \(E \times^{G} F\) . If \(\pi: E \times F \to E \times^{G} F\) denotes the canonical surjection, \(\pi^{-1}(U)\) is an open and closed subset of \(E \times F\) which is stable under G. Since E is connected, there exists a subset \(F' \subset F\), stable under G, such that \(\pi^{-1}(U) = E \times F'\), whence \(U = \pi(E \times F')\). Let \(F'\) and \(F''\) be subsets of F stable under G such that \(\pi(E \times F') = \pi(E \times F'')\). Since E is not empty, we have \(F' = F''\). The map \(F' \mapsto \pi(E \times F')\) is a bijection from the set of G-stable subsets of F onto the set of clopen subsets of \(E \times^{G} F\). The proposition follows.

PROPOSITION 10. --- Let B be a topological space, let (E,p) be a principal covering of B with group G, and let X be a covering of B. Suppose that E and X are connected and nonempty. The following properties are equivalent:

\begin{enumerate}
\def\labelenumi{(\roman{enumi})}
\item
  The covering X is associated to the principal covering E;
\item
  There exists a subgroup H of G such that X is B-isomorphic to E/H;
\item
  There exists a surjective B-morphism \(h: \mathrm{E} \to \mathrm{X}\) ;
\item
  For every point \((y, x)\) of \(E \times_{B} X\) there exists a B-morphism \(r: E \to X\) such that \(r(y) = x\) .
\end{enumerate}

Suppose these conditions are satisfied, and let H be a subgroup of G such that X is B-isomorphic to E/H. The covering X is Galois if and only if the subgroup H is normal in G.

(i)⇒(ii): Let F be a discrete G-set such that the covering \(E \times^{G} F\) is B-isomorphic to X. The set F is nonempty and the space \(E \times^{G} F\) is connected, hence the group G acts transitively on F (prop. 9). Then, the space \(E \times^{G} F\) is B-isomorphic to E/H, where H is the subgroup of G fixing a point of F (I, p. 106, example 4).

(ii)⇒(iii): Indeed, the canonical surjection from E onto E/H is a surjective B-morphism.

(iii)⇒(iv): Let \((y, x)\) be a point of \(E \times_{B} X\). Since the map h is surjective, there exists a point \(y'\) of E such that \(h(y') = x\). The points y and \(y'\) have the same projection in B. Since the covering E is principal with group G, there exists \(g \in G\) such that \(y \cdot g = y'\). The map \(r: E \to X\) defined by \(z \mapsto h(z \cdot g)\) is a B-morphism and one has \(r(y) = x\).

(iv)⇒(i): Since X is nonempty and the projection of E onto B is surjective, the space \(E \times_{B} X\) is not empty. There therefore exists a B-morphism r from E to X. The map ( \(Id_{E}, r\) ): \(E \to E \times_{B} X\) is a section of the covering \(p^{*}(X)\) of E. Since the space E is connected, it follows from Corollary 2 of Prop. 1 in I, p. 69 that the covering \(p^{*}(X)\) is trivializable, which proves that the covering X is associable to the Galois covering p (Proposition 8).

Suppose now that these conditions are satisfied. Let X denote the B-space E/H, q its projection, and \(h: E \rightarrow E/H\) the canonical map.

If the group H is normal in G, then X is a principal covering with group G/H (I, p. 106, example 4). Since X and B are connected and nonempty, X is a Galois covering of B (I, p. 102, th. 2). Conversely, suppose that E/H is a Galois covering of B and let K = Aut \(_{B}\) (E/H) \(^{\circ}\) . We define a map \(\alpha: \mathrm{E} \times \mathrm{G} \to \mathrm{K}\) by associating with \((t, g) \in \mathrm{E} \times \mathrm{G}\) the unique element \(k\) of K such that \(h(t \cdot g) = h(t) \cdot k\) . It is continuous because it is obtained by composing the continuous map \((t, g) \mapsto (h(t), h(t \cdot g))\) from E \(\times\) G into X \(\times_{\mathrm{B}}\) X with the canonical trivialization (I, p. 95, remark 1) X \(\times_{\mathrm{B}}\) X \(\to\) X \(\times\) K of \(q^{*}(\mathrm{X})\) . Since E is connected and K is a discrete group, the continuous map \(t \mapsto \alpha(t, g)\) is constant for every \(g \in \mathrm{G}\); we denote its value by \(\alpha(g)\). If \(t \in \mathrm{E}\), \(g \in \mathrm{G}\), and \(g' \in \mathrm{H}\), we then have

\[
h (t) = h \left(t \cdot g ^ {- 1} g\right) = h \left(t \cdot g ^ {- 1}\right) \cdot \alpha (g) = h \left(t \cdot g ^ {- 1} \cdot g ^ {\prime}\right) \cdot \alpha (g) = h \left(t \cdot \left(g ^ {- 1} g ^ {\prime} g\right)\right).
\]

It follows that \(g^{-1}g'g\) belongs to H, and therefore H is normal in G.

COROLLARY. --- Let E be a Galois covering of B and let X be a covering of B, connected and nonempty. Suppose that X is a finite covering or that the space B is locally connected. If there exists a B-morphism h: E → X, then the covering X is associated with the principal covering E.

In both cases, the B-morphism h is a covering (I, p. 77, cor. 3 of th. 1, and I, p. 78, prop. 5), and a nonempty covering of a connected space is surjective (I, p. 68).

THEOREM 3. --- Let B be a nonempty, connected, and locally connected topological space. Every covering of B can be associated with a Galois covering; every finite covering of B can be associated with a finite Galois covering.

Let X be a covering of B, and denote its projection by q. Since the space B is connected and nonempty, the covering X has a degree; denote this degree by F and endow the set F with the discrete topology. Since the space B is locally connected, the sheaf \(\mathcal{I} = \underline{\mathcal{I}som}_{\mathrm{B}}(\mathrm{B} \times \mathrm{F}; \mathrm{X})\) is locally constant, and at every point b of B, its stalk \(I_{b}\) is canonically isomorphic to the set \(\mathcal{B}(\mathrm{F}; \mathrm{X}_{b})\) of bijections from F onto the fiber \(X_{b}\) (I, p. 89, prop. 12). Let \(E = E_{\mathcal{I}}\) be the associated covering (I, p. 86, corollary).

Let G be the group of permutations of F. For every \(g \in G\), we define a morphism \(\gamma'(g)\) of the sheaf \(\mathcal{I}\) into itself as follows: for every open set U of B, the map \(\gamma'(g)_{\mathrm{U}}\) associates to a U-isomorphism \(\varphi: \mathrm{U} \times \mathrm{F} \to q^{-1}(\mathrm{U})\) the U-isomorphism defined by \((b, f) \mapsto \varphi(b, g(f))\) for \(b \in U\) and \(f \in F\). We have \(\gamma'(\mathrm{Id}_{\mathrm{F}}) = \mathrm{Id}_{\mathcal{I}}\) and \(\gamma'(g \circ g') = \gamma'(g') \circ \gamma'(g)\), for \(g, g' \in G\), so that for every \(g \in G\), \(\gamma'(g)\) is an automorphism of the sheaf \(\mathcal{I}\). The B-morphism \(\gamma(g): E \to E\) associated with \(\gamma'(g)\) is an automorphism (I, p. 50). If \(x = [\mathrm{U}, \varphi, b]\) is an element of E, where U is an open subset of B, b a point of U, and \(\varphi\) a U-isomorphism from U×F to \(q^{-1}(\mathrm{U})\), we have \(\gamma(g)(x) = [\mathrm{U}, \psi, b]\), where \(\psi\) is defined by \(\psi(a, f) = \varphi(a, g(f))\), for \(a \in U\) and \(f \in F\). If g and \(g'\) are elements of G, we have \(\gamma(g \circ g') = \gamma(g') \circ \gamma(g)\). We have \(\gamma(\mathrm{Id}_{\mathrm{F}}) = \mathrm{Id}_{\mathrm{E}}\). We thus define a continuous right action of G on E by setting \(x \cdot g = \gamma(g)(x)\), for \(x \in E\) and \(g \in G\). The group G acts simply transitively on every fiber of E, so that the covering E equipped with this action is a principal covering with group G (I, p. 97, prop. 4). Let h be the map from E × F to X defined by \(h([\mathrm{U}, \varphi, b], f) = \varphi(b, f)\) for every open subset U of B, every point b of U, and every U-isomorphism \(\varphi\) from U × F onto \(q^{-1}(\mathrm{U})\). By definition of the topology on E, it is continuous; it is a B-morphism. For every element g of G and every point \((x, f)\) of E × F, we have \(h(x, g(f)) = h(x \cdot g, f)\). The map h therefore defines, by passing to the quotient, a B-morphism \(h': E \times^{G} F \to X\). For every point \(b\) of B, the map \(h_b\colon \mathrm{E}_b\times \mathrm{F}\to \mathrm{X}_b\) is identified with the map \((\varphi, f)\mapsto \varphi(f)\) from \(\mathcal{B}(\mathrm{F};\mathrm{X}_b)\times \mathrm{F}\) into \(\mathrm{X}_b\), so that the map \(h_b^\prime\) is bijective. Consequently, \(h^\prime\) is a B-isomorphism from \(\mathrm{E}\times^{\mathrm{G}}\mathrm{F}\) onto X (I, p. 30, cor. 2 of prop. 6).

Since the space B is connected, locally connected, and nonempty, the covering X, associable to the principal covering E, is associable to a Galois covering (I, p. 110, corollary 2). If the covering X is finite, then so is the covering E, hence X is associable to a finite Galois covering (loc. cit.).

\subsection{Principal fiber bundles defined by cocycles}

DEFINITION 6. --- Let B be a topological space, let G be a topological group, and let \(\mathcal{U} = (\mathrm{U}_{i})_{i\in \mathrm{I}}\) be an open covering of B. A continuous 1-cocycle on B with values in G, subordinate to \(\mathcal{U}\), is the data, for every pair \((i,j)\) of elements of I, of a continuous map \(g_{i,j}\) from \(\mathrm{U}_i\cap \mathrm{U}_j\) into G, such that for every triple \((i,j,k)\in \mathrm{I}\times \mathrm{I}\times \mathrm{I}\) and every point b of \(\mathrm{U}_i\cap \mathrm{U}_j\cap \mathrm{U}_k\), one has

\[
g _ {i, k} (b) = g _ {i, j} (b) g _ {j, k} (b).
\]

We denote \(Z_{\mathrm{cont}}^{1}(\mathrm{B},\mathcal{U},\mathrm{G})\) the set of continuous 1-cocycles on B with values in G, subordinate to the covering U.

In this issue, we will simply say cocycle instead of continuous 1-cocycle.

Let \((\mathrm{E},p)\) be a principal G-bundle over B and let \(\mathcal{U} = (\mathrm{U}_{i})_{i\in\mathrm{I}}\) be a covering of B by open sets \(U_{i}\). We say that E is trivializable over U if, for every \(i\in I\), E is trivializable over \(U_{i}\). We then call a U-trivialization of E a family \((f_{i})_{i\in\mathrm{I}}\), where \(f_{i}\colon p^{-1}(\mathrm{U}_{i})\to\mathrm{U}_{i}\times\mathrm{G}\) is a trivialization of E over \(U_{i}\).

Let \(\mathcal{U} = (\mathrm{U}_i)_{i\in \mathrm{I}}\) be a covering of B by open sets, and let \((\mathrm{E},p)\) be a principal G-bundle over B equipped with a \(\mathcal{U}\)-trivialization \((f_i)_{i\in \mathrm{I}}\). For every pair \((i,j)\in \mathrm{I}\times \mathrm{I}\), denote by \(g_{ij}\) the map from \(\mathrm{U}_i\cap \mathrm{U}_j\) into G defined by

\[
(b, g _ {i j} (b)) = f _ {i} \circ f _ {j} ^ {- 1} (b, e),
\]

where e denotes the identity element of G. It is continuous.

Since \(f_{i}\) and \(f_{j}\) are compatible with the operations of G, we have

\[
(f _ {i} \circ f _ {j} ^ {- 1}) (b, g) = (b, g _ {i j} (b) g)
\]

for \(b \in \mathrm{U}_i \cap \mathrm{U}_j\) and \(g \in \mathrm{G}\). If \(b\) is a point of \(\mathrm{U}_i \cap \mathrm{U}_j \cap \mathrm{U}_k\) (where \(i, j, k \in \mathrm{I}\)), we have

\[
(f _ {i} \circ f _ {k} ^ {- 1}) (b, e) = (b, g _ {i k} (b))
\]

and

\[
\begin{array}{c} (f _ {i} \circ f _ {k} ^ {- 1}) (b, e) = (f _ {i} \circ f _ {j} ^ {- 1}) \circ (f _ {j} \circ f _ {k} ^ {- 1}) (b, e) \\ = (f _ {i} \circ f _ {j} ^ {- 1}) (b, g _ {j k} (b)) = (b, g _ {i j} (b) g _ {j k} (b)). \end{array}
\]

It follows that

\[
g _ {i k} (b) = g _ {i j} (b) g _ {j k} (b),
\]

so that the family \((g_{ij})\) is a cocycle on B with values in G, subordinate to the covering U. It is called the cocycle defined by the family of trivializations \((f_{i})_{i\in\mathrm{I}}\).

Let B, G, and U be as in Definition 6, and let \((g_{ij}) \in Z_{\mathrm{cont}}^{1}(\mathrm{B}, \mathcal{U}, \mathrm{G})\) be a cocycle. For every pair \((i, j) \in \mathrm{I} \times \mathrm{I}\), the map \(\gamma_{ij} \colon (\mathrm{U}_{i} \cap \mathrm{U}_{j}) \times \mathrm{G} \to (\mathrm{U}_{i} \cap \mathrm{U}_{j}) \times \mathrm{G}\) defined by

\[
\gamma_ {i j} (b, g) = \left(b, g _ {i j} (b) g\right) \quad \text {   for   } b \in \mathrm{U} _ {i} \cap \mathrm{U} _ {j} \text {   and   } g \in \mathrm{G}\tag{5}
\]

is an isomorphism of principal bundles over the base \(U_{i} \cap U_{j}\) . For every triple \((i,j,k) \in \mathrm{I} \times \mathrm{I} \times \mathrm{I}\) and every pair \((b,g) \in (\mathrm{U}_{i} \cap \mathrm{U}_{j} \cap \mathrm{U}_{k}) \times \mathrm{G}\) , we have:

\[
\gamma_ {i k} (b, g) = \gamma_ {i j} \circ \gamma_ {j k} (b, g).
\]

Let F be the topological space obtained by gluing together the spaces \(U_{i} \times G\) along the \((\mathrm{U}_{i} \cap \mathrm{U}_{j}) \times \mathrm{G}\) by means of the bijections \(\gamma_{ij}\) (TG, I, p. 16). For every \(i \in I\), the image of \(U_{i} \times G\) in F is an open set of F (TG, I, p. 17, prop. 9). The canonical projections \(p_{i}: U_{i} \times G \to U_{i}\) glue together into a continuous map p from F to B. The maps \(\gamma_{ij}\) being compatible with the right G-operations in the spaces \(U_{i} \times G\), one deduces a continuous right action of G on F which makes F a principal fiber bundle with base B and group G, equipped with a trivialization over each \(U_{i}\). We say that F is the principal fiber bundle defined by the cocycle \((g_{ij})\).

DEFINITION 7. --- Let B be a topological space, let G be a topological group, and let \(\mathcal{U} = (\mathrm{U}_{i})_{i\in \mathrm{I}}\) be an open covering of B. We say that two cocycles \((g_{ij})\) and \((g_{ij}^{\prime})\) in \(\mathrm{Z}_{\mathrm{cont}}^{1}(\mathrm{B},\mathcal{U},\mathrm{G})\) are cohomologous if there exists a family \((h_{i})_{i\in\mathrm{I}}\) of continuous maps \(h_{i}\colon U_{i}\to G\) such that one has

\[
g _ {i j} ^ {\prime} (b) = h _ {i} (b) g _ {i j} (b) h _ {j} (b) ^ {- 1}\tag{6}
\]

for every pair \((i,j)\in \mathrm{I}\times \mathrm{I}\) and every \(b\in \mathrm{U}_i\cap \mathrm{U}_j\).

The relation “\((g_{ij})\) is cohomologous to \((g'_{ij})\)” is an equivalence relation on the set \(Z_{\mathrm{cont}}^{1}(\mathrm{B},\mathcal{U},\mathrm{G})\). We denote by \(H_{\mathrm{cont}}^{1}(\mathrm{B},\mathcal{U},\mathrm{G})\) the quotient set of \(Z_{\mathrm{cont}}^{1}(\mathrm{B},\mathcal{U},\mathrm{G})\) by this equivalence relation.

PROPOSITION 11. --- Let B be a topological space, G a topological group, and \(\mathcal{U} = (\mathrm{U}_{i})_{i\in \mathrm{I}}\) an open covering of B.

\begin{enumerate}
\def\labelenumi{\alph{enumi})}
\item
  Every principal G-bundle over B that is trivializable over U is isomorphic to a principal bundle defined by a cocycle in \(Z_{\mathrm{cont}}^{1}(\mathrm{B}, \mathcal{U}, \mathrm{G})\).
\item
  Let (E, p) and (E', p') be principal G-bundles over B that are trivializable over U. Let \((f_i)_{i \in I}\) (resp. \((f_i')_{i \in I}\)) be a trivialization of (E, p) (resp. of \((E', p')\)) adapted to U, and denote by \((g_{ij})_{(i,j) \in I \times I}\) (resp. \((g_{i,j}')_{(i,j) \in I \times I}\)) the cocycle defined by this trivialization. Then, the principal fiber bundles (E, p) and (E', p') are isomorphic if and only if these cocycles are cohomologous.
\end{enumerate}

Let us prove \(b\). Let \(\varphi: \mathrm{E} \to \mathrm{E}'\) be an isomorphism of principal bundles with base B and group G. For every \(i \in \mathrm{I}\), let \(h_i\) be the continuous map from \(\mathrm{U}_i\) into G defined by

\[
(b, h _ {i} (b)) = f _ {i} ^ {\prime} \circ \varphi \circ f _ {i} ^ {- 1} (b, e).
\]

Since, for every \(i \in I\) , \(f_{i}\) and \(f_{i}^{\prime}\) are compatible with the operations of G, we have for every \(b \in U_{i}\) and every \(g \in G\) ,

\[
(f _ {i} ^ {\prime} \circ \varphi \circ f _ {i} ^ {- 1}) (b, g) = (b, h _ {i} (b) g)
\]

and

\[
(f _ {i} \circ \varphi^ {- 1} \circ (f _ {i} ^ {\prime}) ^ {- 1}) (b, g) = (b, h _ {i} (b) ^ {- 1} g).
\]

Consequently, for every pair \((i,j)\in \mathrm{I}\times \mathrm{I}\) and every point \(b\) of \(\mathrm{U}_i\cap \mathrm{U}_j\),

\[
\begin{array}{c} f _ {i} ^ {\prime} \circ (f _ {j} ^ {\prime}) ^ {- 1} (b, e) = (f _ {i} ^ {\prime} \circ \varphi \circ f _ {i} ^ {- 1}) \circ (f _ {i} \circ f _ {j} ^ {- 1}) \circ (f _ {j} \circ \varphi^ {- 1} \circ (f _ {j} ^ {\prime}) ^ {- 1}) (b, e) \\ = (b, h _ {i} (b) g _ {i j} (b) h _ {j} (b) ^ {- 1}) \end{array}
\]

so that we have

\[
g _ {i j} ^ {\prime} (b) = h _ {i} (b) g _ {i j} (b) h _ {j} (b) ^ {- 1}.
\]

This proves that the cocycles \((g_{ij})\) and \((g'_{ij})\) are cohomologous.

Conversely, suppose that these cocycles are cohomologous, and let \((h_{i})_{i\in\mathrm{I}}\) be a family of continuous maps \(h_{i}\colon U_{i}\to G\) such that one has \(g_{ij}^{\prime}(b)=h_{i}(b)g_{ij}(b)h_{j}(b)^{-1}\) for \(i,j\in I\) and \(b\in U_{i}\cap U_{j}\). For \(i\in I\), let \(\varphi_{i}\colon p^{-1}(U_{i})\to(p')^{-1}(U_{i})\) be the map defined by \(f_{i}^{\prime}\circ\varphi_{i}\circ f_{i}^{-1}(b,g)=(b,h_{i}(b)g)\), for \(b\in U_{i}\) and \(g\in G\). This is an isomorphism of principal fiber bundles with base \(U_{i}\) and group G. For \((i,j)\in I\times I\), denote by \(\gamma_{ij}\) and \(\gamma_{ij}^{\prime}\) the gluing homeomorphisms associated as above to the cocycles \((g_{ij})\) and \((g_{ij}^{\prime})\), respectively (I, p. 115, formula (5)), so that \(f_{i}(b,g)=\gamma_{ij}\circ f_{j}(b,g)\) and \(f_{i}^{\prime}(b,g)=\gamma_{ij}^{\prime}\circ f_{j}^{\prime}(b,g)\) for all \((b,g)\in(U_{i}\cap U_{j})\times G\). Consequently, for \((i,j)\in I\times I\) and \((b,g)\in(U_{i}\cap U_{j})\times G\), we have the relations

\[
\begin{array}{r l} & f _ {i} ^ {\prime} \circ \varphi_ {j} \circ f _ {i} ^ {- 1} (b, g) = \gamma_ {i j} ^ {\prime} \circ (f _ {j} ^ {\prime} \circ \varphi_ {j} \circ f _ {j} ^ {- 1}) (b, g _ {i j} (b) ^ {- 1} g) \\ & \qquad = \gamma_ {i j} ^ {\prime} (b, h _ {j} (b) g _ {i j} (b) ^ {- 1} g) \\ & \qquad = (b, g _ {i j} ^ {\prime} (b) h _ {j} (b) g _ {i j} (b) ^ {- 1} g) \\ & \qquad = (b, h _ {i} (b) g) = f _ {i} ^ {\prime} \circ \varphi_ {i} \circ f _ {i} ^ {- 1} (b, g). \end{array}
\]

This demonstrates that \(\varphi_{i}\) and \(\varphi_{j}\) coincide on \(p^{-1}(\mathrm{U}_{i}\cap\mathrm{U}_{j})\). The morphisms \(\varphi_{i}\) therefore glue together into a B-morphism of principal bundles from E to E'. The assertion \(b\) follows because any morphism of principal bundles with base B and group G is an isomorphism (I, p. 93, prop. 1).

Now let us prove \(a\). Let \((\mathrm{E}, p)\) be a principal fiber bundle with structure group G, equipped, for every \(i \in \mathrm{I}\), with a trivialization \(f_i: p^{-1}(\mathrm{U}_i) \to \mathrm{U}_i \times \mathrm{G}\) over \(\mathrm{U}_i\). Let \((g_{ij})\) be the cocycle defined by this family, and let F be the principal fiber bundle defined by the cocycle \((g_{ij})\). By construction, the principal fiber bundle F is equipped with a trivialization over \(\mathcal{U}\); this trivialization defines the cocycle \((g_{ij})\). By assertion \(b\), the principal fiber bundle \((\mathrm{E}, p)\) is isomorphic to F.

Let B be a topological space and G a topological group. Let \((E,p)\) be a principal fiber bundle with base B and group G. Let s be a section of the surjective map p (E, II, p. 19, prop. 8). The map \(f\colon B\times G\to E\) defined by \(f(b,g)=s(b)\cdot g\) is a bijection compatible with the action of G. Equip \(B\times G\) with the topology obtained by transport of structure; \((B\times G,\mathrm{pr}_{1})\) is then a principal fiber bundle with base B and group G isomorphic to E. There therefore exists a set T of principal bundles with base B and group G such that every principal fiber bundle with base B and group G is isomorphic to an element of T. Let P(B, G) denote the set of isomorphism classes of principal bundles with base B and group G (E, II, p. 47).

Let \(\mathcal{U} = (\mathrm{U}_i)_{i\in \mathrm{I}}\) be an open cover of B. Let us denote by \(\mathrm{P(B,\mathcal{U},G)}\) the subset of \(\mathrm{P(B,G)}\) consisting of the isomorphism classes of principal bundles that are trivializable over \(\mathcal{U}\). By Proposition 11, there exists a map \(r_{\mathcal{U}}\) from \(\mathrm{H}_{\mathrm{cont}}^{1}(\mathrm{B},\mathcal{U},\mathrm{G})\) into \(\mathrm{P(B,\mathcal{U},G)}\) which associates with the class of a cocycle \((g_{ij})\in Z_{\mathrm{cont}}^{1}(\mathrm{B},\mathcal{U},\mathrm{G})\) the isomorphism class of the principal bundle defined by this cocycle; it is a bijection (loc. cit.). The inverse bijection \(s_{\mathcal{U}}\) associates with the isomorphism class in \(\mathrm{P(B,\mathcal{U},G)}\) of a principal fiber bundle E, trivializable over \(\mathcal{U}\), the class of the cocycle defined by any family of trivializations of E over \(\mathcal{U}\). In this correspondence, the isomorphism class of the trivial principal fiber bundle \(\mathrm{B}\times \mathrm{G}\) corresponds to the cohomology class of the constant cocycle \(g_{ij} = e\), also called the trivial cocycle.

By virtue of the definition of a principal fiber bundle, the set P(B, G) is the union of the sets of the form P(B, U, G), where U is an open covering of B.

Let \(\mathcal{U} = (\mathrm{U}_i)_{i\in \mathrm{I}}\) be an open cover of B and let \(\mathcal{V} = (\mathrm{V}_k)_{k\in \mathrm{K}}\) be an open cover of B finer than \(\mathcal{U}\). We have \(\mathrm{P(B,\mathcal{U},G)}\subset\) \(\mathrm{P(B,\mathcal{V},G)}\); we denote by \(i_{\mathcal{V}\mathcal{U}}\) the canonical injection defined by this inclusion. Let us choose a map \(\varphi \colon \mathrm{K}\to \mathrm{I}\) such that \(\mathrm{V}_k\subset \mathrm{U}_{\varphi (k)}\) for all \(k\in \mathrm{K}\). Given a cocycle \((g_{ij})\in \mathrm{Z}_{\mathrm{cont}}^1 (\mathrm{B},\mathcal{U},\mathrm{G})\), set, for every pair \((k,\ell)\in \mathrm{K}\times \mathrm{K}\), \(\overline{g}_{k\ell} = g_{\varphi (k)\varphi (\ell)}\mid \mathrm{V}_k\cap \mathrm{V}_\ell\).

The family \((\overline{g}_{k\ell})\) is a cocycle on B, with values in G, subordinate to V. If \((g'_{ij}) \in \mathrm{Z}_{\mathrm{cont}}^{1}(\mathrm{B}, \mathcal{U}, \mathrm{G})\) is a cocycle cohomologous to the cocycle \((g_{ij})\), the cocycle \((\overline{g}'_{k\ell})\) deduced from \((g'_{ij})\) is cohomologous to the cocycle \((\overline{g}_{k\ell})\). This yields a map

\[
c (\varphi) \colon \mathrm{H} _ {\mathrm{cont}} ^ {1} (\mathrm{B}, \mathcal {U}, \mathrm{G}) \to \mathrm{H} _ {\mathrm{cont}} ^ {1} (\mathrm{B}, \mathcal {V}, \mathrm{G})
\]

which to the class of \((g_{ij})\) associates the class of \((\overline{g}_{k\ell})\) .

Let E be a principal fiber bundle with base B and group G and, for every \(i \in I\), let \(f_i: E_{U_i} \to U_i \times G\) be a trivialization of the \(U_i\)-principal fiber bundle \(E_{U_i}\). For all \(k \in K\), let \(f_k' : E_{V_k} \to V_k \times G\) be the trivialization deduced from \(f_{\varphi(k)}\) by passing to subsets. Let \((g_{ij}) \in Z_{\mathrm{cont}}^1(\mathrm{B}, \mathcal{U}, \mathrm{G})\) be the cocycle defined by the family \((f_i)\); the cocycle defined by the family \((f_{k}^{\prime})\) is precisely the cocycle \((\overline{g}_{k\ell})\) defined above. Thus, the following diagram is commutative:

\[
\begin{array}{c} \mathrm{H} _ {\mathrm{cont}} ^ {1} (\mathrm{B}, \mathcal {U}, \mathrm{G}) ^ {c (\varphi)} \longrightarrow \mathrm{H} _ {\mathrm{cont}} ^ {1} (\mathrm{B}, \mathcal {V}, \mathrm{G}) \\ \Big \downarrow^ {r _ {\mathcal {U}}} \qquad \qquad \qquad \qquad \Big \downarrow^ {r _ {\mathcal {V}}} \\ \mathrm{P} (\mathrm{B}, \mathcal {U}, \mathrm{G}) ^ {i _ {\mathcal {V} \mathcal {U}}} \longrightarrow \mathrm{P} (\mathrm{B}, \mathcal {V}, \mathrm{G}). \end{array}
\]

The maps \(r_{\mathcal{U}}\) and \(r_{\mathcal{V}}\) being bijective, the map \(c(\varphi)\), just as \(i_{\mathcal{V}\mathcal{U}}\), is an injection and does not depend on the choice of \(\varphi\). Henceforth we shall write \(c_{\mathcal{V}\mathcal{U}}\) instead of \(c(\varphi)\).

Let \(\mathcal{R}\) be the set of elements of \(\mathfrak{P}(\mathfrak{P}(B))\) that are open coverings of B. For every open covering \(\mathcal{U}\) of B, there exists an open covering \(\mathcal{V}\) belonging to \(\mathcal{R}\) such that \(\mathcal{V}\) is both finer and coarser than \(\mathcal{U}\). The set \(\mathcal{R}\) is ordered and filtering for the relation \(\leqslant\) defined by \(\mathcal{U} \leqslant \mathcal{V}\) if \(\mathcal{V}\) is a finer covering than \(\mathcal{U}\). It follows from the preceding that one has defined an inductive system \((\mathrm{H}_{\mathrm{cont}}^{1}(\mathrm{B},\mathcal{U},\mathrm{G}),c_{\mathcal{V}\mathcal{U}})\) relative to the filtered ordered set \(\mathcal{R}\) and that the family \((r_{\mathcal{U}})\) is an inductive system of bijective maps from \((\mathrm{H}_{\mathrm{cont}}^{1}(\mathrm{B},\mathcal{U},\mathrm{G}),c_{\mathcal{V}\mathcal{U}})\) into \((\mathrm{P}(\mathrm{B},\mathcal{U},\mathrm{G}),i_{\mathcal{V}\mathcal{U}})\). If we denote by \(\mathrm{H}_{\mathrm{cont}}^{1}(\mathrm{B},\mathrm{G})\) the inductive limit of the system \((\mathrm{H}_{\mathrm{cont}}^{1}(\mathrm{B},\mathcal{U},\mathrm{G}),c_{\mathcal{V}\mathcal{U}})\) and \(r\colon \mathrm{H}_{\mathrm{cont}}^{1}(\mathrm{B},\mathrm{G})\to \mathrm{P}(\mathrm{B},\mathrm{G})\) the inductive limit of the family \((r_{\mathcal{U}})\), we therefore have:

THEOREM 4. --- The map \(r\colon \mathrm{H}_{\mathrm{cont}}^{1}(\mathrm{B},\mathrm{G})\to \mathrm{P}(\mathrm{B},\mathrm{G})\) is bijective.

Let \(\mathcal{U} = (\mathrm{U}_i)_{i\in \mathrm{I}}\) be an open covering of B; let us denote by \(c_{\mathcal{U}}\) the canonical map \(\mathrm{H}_{\mathrm{cont}}^{1}(\mathrm{B},\mathcal{U},\mathrm{G})\to \mathrm{H}_{\mathrm{cont}}^{1}(\mathrm{B},\mathrm{G})\). If \((g_{ij})\) is a cocycle on B, with values in G, subordinate to the open covering \(\mathcal{U}\), the element \(c_{\mathcal{U}}((g_{ij}))\) of \(\mathrm{H}_{\mathrm{cont}}^{1}(\mathrm{B},\mathrm{G})\) is called the cohomology class of the cocycle \((g_{ij})\).

\section{SIMPLY CONNECTED SPACES}

\subsection{Universal covering}

DEFINITION 1. --- A pointed set is a set X equipped with one of its elements, called the base point. The set X equipped with the point x is sometimes denoted (X, x). Let (X, x) and (Y, y) be pointed sets; a pointed map from (X, x) to (Y, y) is a map f from X to Y such that \(f(x) = y\) .

The notions of pointed topological space, pointed continuous map, pointed covering of a pointed topological space, etc., are defined analogously.

If \((X,x)\) and \((Y,y)\) are pointed topological spaces, the set of pointed continuous maps from \((X,x)\) to \((Y,y)\) is denoted \(\mathcal{C}((X,x);(Y,y))\).

Instead of saying “let f be a pointed map from \((X, x)\) on \((Y, y)\) ”, one often uses the following phrase: “let \(f: (X, x) \to (Y, y)\) be a pointed map”.

Let B be a topological space and let b be a point of B.

DEFINITION 2. --- A pointed covering (E,x) of (B,b) is said to be a universal covering if, for every pointed covering (E',x') of (B,b), there exists a unique morphism of coverings of B, \(f: \mathrm{E} \to \mathrm{E}'\) , such that \(f(x) = x'\) .

If (E, x) and (E', x') are universal coverings of (B, b), the unique B-morphism from (E, x) to (E', x') is an isomorphism of B-spaces.

Let E be a connected covering of B and let x be a point of the fibre \(E_{b}\) . Suppose that, for every pointed covering \((\mathrm{E}^{\prime}, x^{\prime})\) of \((\mathrm{B}, b)\) , there exists a morphism \(f: E \to E^{\prime}\) of coverings of B such that \(f(x) = x^{\prime}\) . Such a morphism f is then unique (I, p. 34, cor. 1 of prop. 11), hence \((\mathrm{E}, x)\) is a universal covering of \((\mathrm{B}, b)\) . *We shall see later (I, p. 126, cor. of prop. 3) that this is in particular the case if every covering of E is trivializable.*

PROPOSITION 1. --- Let B be a connected and locally connected topological space and let b be a point of B. Let (E, x) be a universal covering of (B, b). Then, E is a Galois covering of B and every covering of B is associated to E.

The space E is locally connected, because B is. Let \(E_{0}\) be the connected component of x in E, so that the space \((\mathrm{E}_{0}, x)\) is a pointed covering of \((\mathrm{B}, b)\) (I, p. 80, cor. 1 of prop. 6). There then exists a unique morphism of coverings \(f: E \to E_{0}\) such that \(f(x) = x\). If i denotes the canonical injection of \(E_{0}\) into E, the map \(i \circ f: E \to E\) is a morphism of coverings that maps x to x, as does the map \(Id_{E}\); since \((\mathrm{E}, x)\) is a universal covering of \((\mathrm{B}, b)\), we therefore have \(i \circ f = Id_{E}\). This implies that i is surjective, hence \(E_{0} = E\). Consequently, E is connected.

Let \(y\) be a point of \(\mathrm{E}_b\) and consider the pointed covering \((\mathrm{E},y)\) of \((\mathrm{B},b)\) ; by hypothesis there exists a unique morphism of coverings \(f\colon \mathrm{E}\to \mathrm{E}\) such that \(f(x) = y\) . The map \(s\colon \mathrm{E}\to \mathrm{E}\times_{\mathrm{B}}\mathrm{E}\) defined by \(t\mapsto (t,f(t))\) is a continuous section of the map \(\mathrm{pr}_1\colon \mathrm{E}\times_{\mathrm{B}}\mathrm{E}\to \mathrm{E}\) . By cor. 4 of I, p. 81, this shows that the covering \(\mathrm{E}\times_{\mathrm{B}}\mathrm{E}\) of E given by the map \(\mathrm{pr}_1\) is trivializable. The covering E is therefore Galois (th. 2 of I, p. 102). It then follows from I, p. 112, corollary of prop. 10, that every covering of B is associated to E.

COROLLARY. --- Let B be a locally connected topological space. Let b be a point of B and let (E, x) be a universal covering of (B, b). For a subspace A of B, the following two properties are equivalent:

\begin{enumerate}
\def\labelenumi{(\roman{enumi})}
\item
  The covering E is trivializable over A;
\item
  Every covering of B is trivializable over A.
\end{enumerate}

Property (ii) obviously implies property (i). The converse follows from the fact that every covering of B is associated to the covering E (I, p. 105, prop. 7).

\subsection{Convex subsets of a numerical space}

Let E be the n-dimensional numerical space \(^{*}\) (or, more generally, a topological vector space over R)*. For every pair \((x,y)\) of points of E, one calls segment (resp. open segment) with endpoints x and y (cf. EVT, II, p. 7) the set of points of E of the form \(tx + (1-t)y\) , for \(t \in [0,1]\) (resp. for \(t \in ]0,1[\) ). Let A be a subset of E. One says that the set A is convex if for every pair \((x,y)\) of points of A and every \(t \in I\) , the point \(tx + (1-t)y\) belongs to A.

*A convex subset is path-connected.*

Lemma 1. --- Let E be the n-dimensional numerical space and let A be a convex and compact subset of E of which 0 is an interior point. For every \(x \in E\), denote by \(p_{\mathrm{A}}(x)\) the lower bound in \(\overline{R}\) of the set of real numbers \(t > 0\) such that \(x \in tA\).

The map \(p_{A}\) is finite, continuous, and satisfies the following properties:

\begin{enumerate}
\def\labelenumi{(\roman{enumi})}
\item
  For every \(x \in E\) such that \(x \neq 0\) , we have \(p_{\mathrm{A}}(x) > 0\) ;
\item
  For every \(s \in R_{+}\) and every \(x \in E\) , we have \(p_{\mathrm{A}}(sx) = sp_{\mathrm{A}}(x)\) .
\item
  For all \(x\) and \(y \in \mathrm{E}\) , we have \(p_{\mathrm{A}}(x + y) \leqslant p_{\mathrm{A}}(x) + p_{\mathrm{A}}(y)\) .
\item
  For a point \(x\) of E to belong to A, it is necessary and sufficient that \(p_{\mathrm{A}}(x) \leqslant 1\) .
\end{enumerate}

For \(x \in E\), denote by \(\|x\|\) the Euclidean norm of x (TG, VI, p. 7). Since A is compact, there exists a real number \(M > 0\) such that every point x of A satisfies \(\|x\| \leqslant M\). Since 0 is an interior point of A, there exists a real number \(m > 0\) such that every point x of E such that \(\|x\| \leqslant m\) belongs to A. Consequently, one has the relation \(\|x\|/\mathrm{M} \leqslant p_{\mathrm{A}}(x) \leqslant \|x\|/m\), for every \(x \in E\). In particular, \(p_{\mathrm{A}}(0) = 0\) and \(p_{\mathrm{A}}(x) \neq 0\) if \(x \neq 0\).

Assertions (ii) and (iv) follow immediately from the definition of the map \(p_{A}\) .

Let \(x\) and \(y\) be points of E. Let \(x'\) and \(y'\) be points of A such that \(x = p_{\mathrm{A}}(x)x'\) and \(y = p_{\mathrm{A}}(y)y'\) . If \(x\) and \(y\) are not both zero, \(p_{\mathrm{A}}(x) + p_{\mathrm{A}}(y) > 0\) and one has

\[
\begin{array}{l} x + y = p _ {\mathrm{A}} (x) x ^ {\prime} + p _ {\mathrm{A}} (y) y ^ {\prime} \\ \qquad = (p _ {\mathrm{A}} (x) + p _ {\mathrm{A}} (y)) \left(\frac {p _ {\mathrm{A}} (x)}{p _ {\mathrm{A}} (x) + p _ {\mathrm{A}} (y)} x ^ {\prime} + \frac {p _ {\mathrm{A}} (y)}{p _ {\mathrm{A}} (x) + p _ {\mathrm{A}} (y)} y ^ {\prime}\right). \end{array}
\]

Since A is convex, this shows that \(x + y\) belongs to \((p_{\mathrm{A}}(x) + p_{\mathrm{A}}(y))\mathrm{A}\) , whence \(p_{\mathrm{A}}(x + y) \leqslant p_{\mathrm{A}}(x) + p_{\mathrm{A}}(y)\) . If x = y = 0, this inequality is still satisfied, because \(p_{\mathrm{A}}(0) = 0\) . This proves assertion (iii).

Applying this inequality to \(x + y\) and -y, one deduces that, for every pair \((x, y)\) of points of E, we have

\[
| p _ {\mathrm{A}} (x + y) - p _ {\mathrm{A}} (x) | \leqslant \max (p _ {\mathrm{A}} (y), p _ {\mathrm{A}} (- y)) \leqslant m ^ {- 1} \| y \|.
\]

This proves that the map \(p_{A}\) is continuous, whence the lemma.

The map \(p_{A}\) is called the gauge of the convex set A.

PROPOSITION 2. --- Let E be the n-dimensional numerical space and let A be a convex and compact subset of E of which 0 is an interior point. There exists a unique bijection u of E onto itself satisfying the following three properties:

\begin{enumerate}
\def\labelenumi{(\roman{enumi})}
\item
  For every \(x \in E\) and every \(t \in R_{+}\) , \(u(tx) = tu(x)\) ;
\item
  For every \(x \in E\) , there exists \(\lambda \in R_{+}\) such that \(u(x) = \lambda x\) ;
\item
  We have \(u(\mathrm{A}) = \mathbf{B}_n\)
\end{enumerate}

The map u is a homeomorphism and induces, by passage to subspaces, a homeomorphism of A onto \(B_{n}\) , a homeomorphism of the interior of A onto the interior of \(B_{n}\) , and a homeomorphism of the boundary of A in E onto the sphere \(S_{n-1}\) .

Let \(p_{A}\) be the gauge of the convex set A. Let \(x \in E\) and let \(t \in R_{+}\); for \(x\) to belong to \(tA\), it is necessary and sufficient that \(p_{\mathrm{A}}(x) \leqslant t\). Since A is compact, there exists a real number M such that \(\|x\| \leqslant M\) for all \(x \in A\).

Let \(u\) be a map satisfying the conditions of the proposition. We have \(u(0) = 0\). Let \(x \in \mathrm{E}\setminus\{0\}\) and let \(\lambda \in \mathbf{R}_+\) such that \(u(x) = \lambda x\). For all \(t \in \mathbf{R}_+\) such that \(tx \in \mathrm{A}\), we have \(u(tx) = t\lambda x\). Since \(u\) is injective, \(\lambda \neq 0\). Since \(u(\mathrm{A})\) is contained in \(\mathbf{B}_n\), we also have \(\lambda \leqslant p_{\mathrm{A}}(x)/\|x\|\). Set \(z = x/\|x\|\); it is a point of \(\mathbf{S}_{n-1}\). For \(z\) to have a preimage in A under \(u\), it is necessary and sufficient that the point \((\lambda \|x\|)^{-1}x\) belong to A, that is, \(\lambda \|x\| \geqslant p_{\mathrm{A}}(x)\), i.e. \(\lambda \geqslant p_{\mathrm{A}}(x)/\|x\|\). This implies the uniqueness of such a map \(u\).

Let us then denote by \(u\) the map of E into itself which sends 0 to 0 and \(x\) onto \((p_{\mathrm{A}}(x) / \| x\|)x\), for every \(x \in \mathrm{E}\setminus\{0\}\).

By lemma 1, \(u\) is continuous at every point of \(\mathrm{E}\setminus\{0\}\). We have \(\| u(x)\| = p_{\mathrm{A}}(x)\) for all \(x\in \mathrm{E}\) and \(p_{\mathrm{A}}(x)\to 0\) when \(x\to 0\) (loc. cit.); consequently, \(u\) is continuous at 0. Hence, \(u\) is continuous.

The only preimage of 0 under u is 0. Let \(y \in E\setminus\{0\}\). For an element x of E to satisfy \(u(x) = y\), it is necessary and sufficient that \(x = (\|y\| / p_{\mathrm{A}}(y))y\). This entails that u is a continuous bijection of E onto itself. Its inverse is the map \(v: y \mapsto (\|y\| / p_{\mathrm{A}}(y))y\). Since \(p_{A}\) is continuous and vanishes only at 0, the map v is continuous at every point of \(E\setminus\{0\}\); the inequality \(p_{\mathrm{A}}(y) \geqslant \|y\| / \mathrm{M}\) entails that \(\|v(y)\| \leqslant \mathrm{M}\|y\|\) for all \(y \in E\). It follows that v is continuous. Hence, the map u is a homeomorphism of E onto itself. Since the relations \(p_{\mathrm{A}}(x) \leqslant 1\) and \(x \in A\) are equivalent, we also have \(u(\mathrm{A}) = \mathbf{B}_{n}\). The proposition follows.

Example. --- The set \([0,1]^{n}\) is a convex, compact subset of \(R^{n}\) with nonempty interior. It follows from proposition 2 of I, p. 123 that it is homeomorphic to the closed Euclidean ball. More precisely, for every point x of the interior \(]0,1[^{n}\) and every point b of the open Euclidean ball, there exists a homeomorphism of \([0,1]^{n}\) onto \(B_{n}\) which sends x to b and induces by passage to subspaces a homeomorphism of \(]0,1[^{n}\) onto the open Euclidean ball, as well as a homeomorphism of the boundary of \([0,1]^{n}\) onto the sphere \(S_{n-1}\).

Consequently, every convex, compact subset of \(R^{n}\) with nonempty interior is homeomorphic to a parallelepiped (loc. cit.).

\subsection{Simply connected spaces}

DEFINITION 3. --- A topological space is said to be simply connected if all its coverings are trivializable.

A simply connected space is connected. Indeed, if a space X is the disjoint union of two nonempty open sets U and V, the canonical injection of U into X is a covering which is not trivializable.

The empty space is simply connected. Every topological space reduced to a single point is simply connected.

Remarks. --- 1) Let B be a simply connected topological space and (E, p) a covering of B. If the space E is connected and nonempty, then p is a homeomorphism. Let, for example, G be a connected topological group and H a discrete subgroup of G, so that the canonical map p: G → G/H makes G a covering of G/H (I, p. 100, cor. 2 of th. 1). If the space G/H is simply connected, then p is a homeomorphism and H is the trivial subgroup.

\begin{enumerate}
\def\labelenumi{\arabic{enumi})}
\setcounter{enumi}{1}
\tightlist
\item
  Let B be a topological space in which every point has a simply connected neighborhood; then every covering of a covering of B is a covering of B. Indeed, consider a covering (E, p) of B and a covering (F, q) of E. Let us prove that (F, \(p \circ q\)) is a covering of B. Since the question is local in B, we may assume that the space B is simply connected and hence that E is a trivializable covering of B. Let V be a connected component of E. It is open and closed and \(p \mid V: V \to B\) is a homeomorphism (I, p. 69, proposition 1); consequently, the space V is simply connected. Every connected component W of \(q^{-1}(V)\) is open and closed in \(q^{-1}(V)\), hence in F, and the map q induces a homeomorphism from W onto V. The map \(p \circ q\) thus makes F a covering of B, trivializable (loc. cit.).
\end{enumerate}

PROPOSITION 3. --- Let B be a topological space. Let (E, p) be a covering of B and let y be a point of E. Let X be a simply connected topological space, let f: X → B be a continuous map and x a point of X such that f(x) = p(y). There exists a unique continuous lift g: X → E of the map f such that g(x) = y.

The maps \(g\) sought correspond bijectively (I, p. 9, prop. 3) to the continuous sections \(s\) of the covering \(\mathrm{pr}_1: \mathrm{X} \times_{\mathrm{B}} \mathrm{E} \to \mathrm{X}\) such that \(s(x) = (x, y)\). Such a section exists because the space X is simply connected; it is unique because the space X is connected (I, p. 34, cor. 1 of prop. 11).

Example 1. --- Let X be a simply connected topological space and let f be a continuous function from X into \(C^{*}\). Recall (I, p. 101, example 6) that the map \(z \mapsto e^{z}\) makes C a covering of \(C^{*}\). By prop. 3, there exists a continuous function \(h: X \to C\) such that \(f(x) = e^{h(x)}\) for all \(x \in X\). If \(h': X \to C\) is another continuous function such that \(f(x) = e^{h'(x)}\) for all \(x \in X\), there exists an integer \(q \in Z\) such that \(h' = h + 2\pi iq\).

Similarly, let \(n\) be an integer \(>0\); the map \(z \mapsto z^{n}\) makes \(C^{*}\) a covering of itself (I, p. 101, example 5). There therefore exists a continuous function \(k\) from X into \(C^{*}\) such that \(k(x)^{n} = f(x)\) for all \(x \in X\), for example the function \(k(x) = e^{h(x)/n}\). If \(k' : X \to C^{*}\) is another continuous function such that \(k'(x)^{n} = f(x)\) for all \(x \in X\), there exists an \(n\)th root of unity \(\mu \in C\) such that \(k' = \mu k\).

COROLLARY. — Let B be a topological space and let b be a point of B. Let E be a simply connected covering of B. For every point x of \(E_{b}\) the pointed space (E, x) is a universal covering of (B, b).

PROPOSITION 4. --- The product of two simply connected spaces, one of which is locally connected, is a simply connected space.

Let X and Y be simply connected spaces, and suppose that Y is locally connected. The space \(X \times Y\) is connected, since X and Y are (I, p. 124). Let \((Z, f)\) be a nonempty covering of \(X \times Y\); let us prove that it is trivializable. By Corollary 2 of I, p. 70, it suffices to prove that for every point \(z_{0}\) of Z, there exists a continuous section s of f such that \(s(f(z_{0})) = z_{0}\). Let therefore \(z_{0}\) be a point of Z. Set \((x_{0}, y_{0}) = f(z_{0})\). The subspace \(X \times \{y_{0}\}\) of X × Y is homeomorphic to X. Consequently, the covering deduced from Z over \(X \times \{y_{0}\}\) is trivializable and has a continuous section σ such that \(\sigma(x_{0}, y_{0}) = z_{0}\). Similarly, for every point x of X, there exists a continuous section \(\tau_{x}\) of f over \(\{x\} \times Y\) such that \(\tau_{x}(x, y_{0}) = \sigma(x, y_{0})\). For \((x, y) \in X \times Y\), set \(s(x, y) = \tau_{x}(x, y)\). The map s is a section of f, and one has \(s(x_{0}, y_{0}) = z_{0}\). Since the space Y is connected and locally connected, it follows from th. 1 of I, p. 35 that the map s is continuous.

PROPOSITION 5. --- A connected and locally connected topological space such that the intersection of any two connected open subsets is connected is a simply connected space.

Let B be such a topological space. Let \((\mathrm{E}, p)\) be a covering of B, let x be a point of E, and denote \(b = p(x)\). It is a matter of proving that there exists a continuous section s of p such that \(s(b) = x\) (cor. 2, I, p. 70). Let S be the set of pairs \((\mathrm{U}, s_{\mathrm{U}})\) where U is a connected open subset of B containing b and \(s_{U}\) is a continuous section of \(p_{\mathrm{U}}: p^{-1}(\mathrm{U}) \to \mathrm{U}\) such that \(s_{\mathrm{U}}(b) = x\). The set S is not empty (I, p. 34, prop. 10). Let \((\mathrm{U}, s_{\mathrm{U}})\) and \((\mathrm{V}, s_{\mathrm{V}})\) be elements of S. Then, \(s_{U} \mid U \cap V\) and \(s_{V} \mid U \cap V\) are continuous sections of \(p_{U \cap V}\) which take the value x at b. By hypothesis, \(U \cap V\) is connected; we therefore have \(s_{U} \mid U \cap V = s_{V} \mid U \cap V\) (I, p. 34, cor. 1 of prop. 11). Let A be the union of the open sets U as \((\mathrm{U}, s_{\mathrm{U}})\) runs through S, and let s: \(A \to E\) be the unique map such that \(s \mid U = s_{U}\) for every pair \((\mathrm{U}, s_{\mathrm{U}}) \in S\). The set A is open and connected (TG, I, p. 81, prop. 2), it contains b, and s is a continuous section of \(p_{A}\) such that \(s(b) = x\). It now suffices to prove that the set A is closed, which will imply that it is equal to B.

Let \(a\) be a point of B adherent to A, and let V be a connected open neighborhood of \(a\) such that the covering E is trivializable over V. There exists a point \(c\) in \(\mathrm{A} \cap \mathrm{V}\) and a continuous section \(s_{\mathrm{V}}\) of \(p_{\mathrm{V}}\) such that \(s_{\mathrm{V}}(c) = s(c)\). Let \(\mathrm{A}'\) be the open set \(\mathrm{A} \cup \mathrm{V}\). Since \(\mathrm{A} \cap \mathrm{V}\) is connected, there exists a continuous section \(s'\) of \(p_{\mathrm{A}'}\) which extends \(s\) and \(s_{\mathrm{V}}\) (I, p. 35, cor. 3 of prop. 11); the pair \((\mathrm{A}', s')\) belongs to \(\mathcal{S}\) and \(\mathrm{A}'\) is therefore contained in A. Consequently, \(a\) belongs to A, and A is closed.

COROLLARY. --- Every interval of the real line R is simply connected.

Indeed, the connected subspaces of R are the intervals (TG, IV, p. 8, th. 4) and the intersection of two intervals is an interval.

Example 2. --- The n-dimensional numerical space \(R^{n}\) (TG, VI, p. 1) is simply connected. Indeed, it is a product of simply connected and locally connected spaces (I, p. 126, prop. 4 and cor. of prop. 5 above). The same holds for every open or closed parallelepiped of \(R^{n}\). A parallelotope and an open or closed Euclidean ball in \(R^{n}\) are simply connected because they are homeomorphic to a block (TG, VI, p. 10, prop. 2 and I, p. 124, example).

PROPOSITION 6. --- Let X be a topological space. Let U₁ and U₂ be open (resp. closed) subspaces of X such that X = U₁ ∪ U₂. Suppose that U₁ and U₂ are simply connected and that their intersection U₁ ∩ U₂ is connected and nonempty. Then, the space X is simply connected.

Let (E, p) be a covering of X and let y be a point of E. It suffices to show that there exists a continuous section s of p such that \(s(p(y)) = y\) (I, p. 70, cor. 2 of prop. 1). Suppose for example that \(p(y)\) belongs to \(U_{1}\). There then exists a continuous section \(s_{1}: U_{1} \to E\) of \(p_{U_{1}}\) such that \(s_{1}(p(y)) = y\). Let x be a point of \(U_{1} \cap U_{2}\); there exists a continuous section \(s_{2}\) of \(p_{U_{2}}\) such that \(s_{2}(x) = s_{1}(x)\). By Corollary 3 of Proposition 11 of I, p. 34, there exists a continuous section s of p extending both \(s_{1}\) and \(s_{2}\).

Examples. --- 3) For \(n \geqslant 2\), the sphere \(\mathbf{S}_{n}\) (TG, VI, p. 10) is simply connected. Indeed, the sphere \(\mathbf{S}_{n}\) is the union of two closed hemispheres homeomorphic to \(\mathbf{B}_{n}\) whose intersection is homeomorphic to

\(S_{n-1}\) (cf. TG, VI, p. 12). For \(n \geqslant 2\), the sphere \(S_{n-1}\) is connected, whence the assertion.

On the other hand, the circle \(S_{1}\) is not simply connected. Indeed, the continuous map \(p: R \to S_{1}\) defined by \(p(\theta) = (\cos \theta, \sin \theta)\) makes R an infinite-degree covering of \(S_{1}\), which is connected and hence not trivializable.

\begin{enumerate}
\def\labelenumi{\arabic{enumi})}
\setcounter{enumi}{3}
\tightlist
\item
  Let E be a finite-dimensional vector space of dimension n over R, and let F be an affine subspace of E of codimension \(p \geqslant 3\). The set \(R^{p}\setminus\{0\}\) is homeomorphic to \(R \times S_{p-1}\) (TG, VI, p. 10, cor. 2), hence it is simply connected, since R and \(S_{p-1}\) are locally connected and simply connected (I, p. 126, prop. 4). The set \(E\setminus F\), homeomorphic to \(R^{n-p} \times (R^{p}\setminus\{0\})\), is therefore simply connected (loc. cit.).
\end{enumerate}

PROPOSITION 7. — Let X be a topological space. Let \(U_{1}\) and \(U_{2}\) be connected open (resp. closed) subspaces of X such that \(X = U_{1} \cup U_{2}\). If the space X is simply connected, then \(U_{1} \cap U_{2}\) is connected.

Let us set \(U = U_{1} \cap U_{2}\) and suppose for contradiction that the space U is not connected. We will construct a connected covering of X of infinite degree; such a covering is not trivializable.

By hypothesis, the set U is the union of two disjoint, nonempty sets A and B that are open (resp. closed) in X. For \(i \in \{1, 2\}\), let \(Y_i\) be the \(U_i\)-space \((\mathrm{U}_i \times \mathbf{Z}, \mathrm{pr}_1)\), and let \(Z_i\) be the subspace \(U \times \mathbf{Z}\) of \(Y_i\). The map \(h: Z_1 \to Z_2\) defined by

\[
h (x, n) = \left\{ \begin{array}{l l} (x, n) & \text {if } x\in A \text { and } n\in \mathbf {Z} \\ (x, n + 1) & \text {if } x\in B \text { and } n\in \mathbf {Z} \end{array} \right.
\]

is a homeomorphism. Let Y be the space obtained by gluing \(Y_{1}\) and \(Y_{2}\) along \(Z_{1}\) and \(Z_{2}\) by means of the homeomorphism h (TG, I, p. 17).

For \(i \in \{1,2\}\), the canonical map \(g_i\) of \(Y_i\) into Y is a homeomorphism of \(Y_i\) onto an open (resp. closed) subset of Y (loc. cit., prop. 9). For every integer \(n \in \mathbf{Z}\), the sets \(g_i(U_i \times \{n\})\), \(i \in \{1,2\}\), are connected; since A and B are nonempty, \(g_1(U_1 \times \{n\})\) meets \(g_2(U_2 \times \{n\})\) and \(g_2(U_2 \times \{n+1\})\). It follows that the space Y is connected (TG, I, p. 81, corollary).

For \(x \in \mathrm{U}\) and \(n \in \mathbf{Z}\), we have \((\mathrm{pr}_1 \circ h)(x, n) = x = \mathrm{pr}_1(x, n)\). There therefore exists a unique continuous map \(p \colon \mathrm{Y} \to \mathrm{X}\) such that \(p \circ g_i = \mathrm{pr}_1\) for \(i \in \{1, 2\}\). Let us prove that the X-space \((\mathrm{Y}, p)\) is a covering.

By construction, the fibers of the map p are homeomorphic to the discrete space Z.

For \(i \in \{1,2\}\), the map \(g_i\) defines, by passing to subspaces, an isomorphism of \(U_i\)-spaces from \(\mathrm{U}_i \times \mathbf{Z}\) onto \(p^{-1}(\mathrm{U}_i)\).

By definition of the space Y, there exists a unique map k from \((\mathrm{X}\setminus\mathrm{A})\times\mathbf{Z}\) into Y such that \(k(x,n)=g_{1}(x,n)\) for \(x\in(\mathrm{U}_{1}\setminus\mathrm{A})\times\mathbf{Z}\) and \(k(x,n)=g_{2}(x,n-1)\) for \(x\in(\mathrm{U}_{2}\setminus\mathrm{A})\times\mathbf{Z}\); it is an isomorphism of \((\mathrm{X}\setminus\mathrm{A})\)-spaces from \((\mathrm{X}\setminus\mathrm{A})\times\mathbf{Z}\) onto \(p^{-1}(\mathrm{X}\setminus\mathrm{A})\). Likewise, there exists a unique map \(k'\) from \((\mathrm{X}\setminus\mathrm{B})\times\mathbf{Z}\) into Y which coincides with \(g_{1}\) on \((\mathrm{U}_{1}\setminus\mathrm{B})\times\mathbf{Z}\) and with \(g_{2}\) on \((\mathrm{U}_{2}\setminus\mathrm{B})\times\mathbf{Z}\); it is an isomorphism of \((\mathrm{X}\setminus\mathrm{B})\)-spaces from \((\mathrm{X}\setminus\mathrm{B})\times\mathbf{Z}\) onto \(p^{-1}(\mathrm{X}\setminus\mathrm{B})\).

This proves that the X-space \((Y,p)\) is trivializable over the parts \(U_{1}\), \(U_{2}\), \(X\setminus A\), and \(X\setminus B\). We have \(U_{1} \cup U_{2} = X\); if \(U_{1}\) and \(U_{2}\) are open in X, this proves that \((Y,p)\) is a covering of X. The same holds when \(U_{1}\) and \(U_{2}\) are closed in X because then \(X\setminus A\) and \(X\setminus B\) are open subsets of X whose union is X. The proposition is thus proved.

COROLLARY. --- Let X be a simply connected topological space and let A be a connected subset of X. If the complement of A is connected, then its boundary is also connected.

Let \(X_{1}\) and \(X_{2}\) be the closures of A and of \(X\setminus A\), respectively. The sets \(X_{1}\) and \(X_{2}\) are closed and connected (TG, I, p. 81, prop. 1); we have \(X_{1} \cup X_{2} = X\) and their intersection \(X_{1} \cap X_{2}\) is equal to the boundary of A. It then suffices to apply Proposition 7.
